\documentclass[compress,dvips,xcolor={dvipsnames},t]{beamer}

\usepackage[T1]{fontenc}
\usepackage[english]{babel}
\usepackage{lmodern}
\usepackage{beramono}          % Bold in verbatim
\usepackage{hyphenat}          % \hyp{} is a breakable dash
\usepackage{xspace}
\usepackage{graphicx}
\usepackage{alltt}
\usepackage{amssymb,amsmath,stmaryrd}
\usepackage{mathpartir}
\usepackage{multicol}
\usepackage{clrscode}
\usepackage{url}
\usepackage{booktabs}

%-*-latex-*-

% Shell prompting
%
\newcommand{\shellprompt}{\textcolor{red}{\texttt{}\$}}
\newcommand{\shellin}[1]{\textsf{\shellprompt{}\ {#1}}}

% OCaml toplevel
%
\newcommand{\camlprompt}{\texttt{\textcolor{blue}{\#}}}
%\newcommand{\semi}{\texttt{\textbf{;;}}}
\newcommand{\topin}[1]{\textsf{\#}\ {#1}\textbf{;;}}
\newcommand{\topout}[1]{\emph{\textsf{\textcolor{blue}{#1}}}}
\newcommand{\error}[1]{\emph{\texttt{\textcolor{red}{#1}}}}

% Tokens
%
\newcommand{\tok}[1]{\small \cst{#1}\xspace}

\newcommand{\Teof}{\tok{EOF}}

\newcommand{\Tident}{\tok{IDENT}}
\newcommand{\Tint}{\tok{INT}}
\newcommand{\Tplus}{\tok{PLUS}}
\newcommand{\Tminus}{\tok{MINUS}}
\newcommand{\Ttimes}{\tok{TIMES}}
\newcommand{\Tslash}{\tok{SLASH}}
\newcommand{\Tequal}{\tok{EQUAL}}
\newcommand{\Tarrow}{\tok{ARROW}}
\newcommand{\Tlpar}{\tok{LPAR}}
\newcommand{\Trpar}{\tok{RPAR}}
\newcommand{\Tlet}{\tok{LET}}
\newcommand{\Trec}{\tok{REC}}
\newcommand{\Tin}{\tok{IN}}
\newcommand{\Tfun}{\tok{FUN}}
\newcommand{\Tmatch}{\tok{MATCH}}
\newcommand{\Tif}{\tok{IF}}
\newcommand{\Tthen}{\tok{THEN}}
\newcommand{\Telse}{\tok{ELSE}}
\newcommand{\Ttrue}{\tok{TRUE}}
\newcommand{\Tfalse}{\tok{FALSE}}
\newcommand{\Tand}{\tok{AND}}
\newcommand{\Tor}{\tok{OR}}
\newcommand{\Tnot}{\tok{NOT}}

% ocamllex
%
\newcommand{\Xrule}{\kwd{rule}\xspace}
\newcommand{\Xparse}{\kwd{parse}\xspace}

% Abstract Syntax
%
%\newcommand{\clos}[3]{\cst{Clos} \, \lpar{#1},#2,#3\rpar}

% Meta-variable
%
\newcommand{\meta}[1]{\overline{#1}}

% Judgements (evaluation and typing)
%
\newcommand{\ceval}[2]{{#1} \twoheadrightarrow {#2}}
\newcommand{\eval}[3]{{#1} \vdash {#2} \twoheadrightarrow {#3}}
\newcommand{\meval}[3]{\eval{#1}{\src{#2}}{#3}}
\newcommand{\evalb}[2]{\ceval{#1}{#2}}
\newcommand{\ieval}[5]{{#1}/{#2} \vdash {#3} \twoheadrightarrow {#4}/{#5}}
\newcommand{\trel}[3]{{#1} \vdash {#2} : {#3}}
\newcommand{\dom}[1]{\mathrm{dom} (#1)}
\newcommand{\codom}[1]{\mathrm{codom} (#1)}
\newcommand{\subst}[3]{{#1}[#2 \leftarrow #3]}

% Syntax analyser
%
\newcommand{\src}[1]{\langle\!\langle{#1}\rangle\!\rangle}
%\newcommand{\src}[1]{\llbracket{#1}\rrbracket}

% Decorated AST node
%
\newcommand{\mknode}[3]%
  {#1~[tnpos=l]{\textbf{#2}}~[tnpos=r]{\textcolor{red}{#3}}}

% Comments and examples
%
\newcommand{\remarque}{\textbf{Remarque}\xspace}
\newcommand{\remarques}{\textbf{Remarques}\xspace}
\newcommand{\exemple}{\textbf{Exemple}\xspace}
\newcommand{\exemples}{\textbf{Exemples}\xspace}

% Regular formal langages
%
%\newcommand{\R}[1]{{\cal R} ({#1}^{*})}
\newcommand{\Reg}[1]{{\cal R} ({#1}^{*})}

% AST annotation (typed expression)
%
\newcommand{\te}[2]{{#1}^{#2}}

% Type Unification
%
\newcommand{\mgu}{\textsf{mgu}\xspace}

%%-*-latex-*-


% OCaml pieces of source code
%
\newcommand{\phrase}[1]{\texttt{#1\texttt{\scriptsize ;;}}}
\newcommand{\excerpt}[1]{\texttt{#1}}

% OCaml language
%
\newcommand{\ipat}[1]{\overline{#1}}
\newcommand{\topclos}{\langle{fun}\rangle}

\newcommand{\wild}{{\LARGE \_}}
\newcommand{\kwd}[1]{\textsf{\textbf{#1}}}
\newcommand{\ident}[1]{\textsf{#1}}
\newcommand{\cst}[1]{\textsf{#1}}
\newcommand{\type}[1]{\textsf{#1}}
\newcommand{\str}[1]{\textsf{"#1"}}
\newcommand{\num}[1]{\textsf{#1}}
\newcommand{\comment}[1]{\textsf{(* #1 *)}}
%\newcommand{\cons}{\textsf{::}}
\newcommand{\semi}{\textsf{;;}}
\newcommand{\lpar}{\textsf{(}}
\newcommand{\rpar}{\textsf{)}}
\newcommand{\lbra}{\textsf{[}}
\newcommand{\rbra}{\textsf{]}}
\newcommand{\equal}{\textsf{=}\xspace}
\newcommand{\nequal}{\textsf{\footnotesize <>}\xspace}
\newcommand{\assign}{\textsf{:=}\xspace}
\newcommand{\unit}{\textsf{()}\xspace}
\newcommand{\vbar}{\texttt{|}\xspace}

\newcommand{\Xopen}{\kwd{open}\xspace}
\newcommand{\Xlet}{\kwd{let}\xspace}
\newcommand{\Xrec}{\kwd{rec}\xspace}
\newcommand{\Xnonrec}{\kwd{nonrec}\xspace}
\newcommand{\Xval}{\kwd{val}\xspace}
\newcommand{\Xin}{\kwd{in}\xspace}
\newcommand{\Xfun}{\kwd{fun}\xspace}
\newcommand{\Xfunction}{\kwd{function}\xspace}
\newcommand{\Xmatch}{\kwd{match}\xspace}
\newcommand{\Xtry}{\kwd{try}\xspace}
\newcommand{\Xwith}{\kwd{with}\xspace}
\newcommand{\Xif}{\kwd{if}\xspace}
\newcommand{\Xthen}{\kwd{then}\xspace}
\newcommand{\Xelse}{\kwd{else}\xspace}
\newcommand{\Xtype}{\kwd{type}\xspace}
\newcommand{\Xand}{\kwd{and}\xspace}
\newcommand{\Xof}{\kwd{of}\xspace}
\newcommand{\Xfor}{\kwd{for}\xspace}
\newcommand{\Xwhile}{\kwd{while}\xspace}
\newcommand{\Xdo}{\kwd{do}\xspace}
\newcommand{\Xdone}{\kwd{done}\xspace}
\newcommand{\Xbegin}{\kwd{begin}\xspace}
\newcommand{\Xend}{\kwd{end}\xspace}
\newcommand{\Xexception}{\kwd{exception}\xspace}
\newcommand{\Xas}{\kwd{as}\xspace}
\newcommand{\Xtrue}{\kwd{true}\xspace}
\newcommand{\Xfalse}{\kwd{false}\xspace}

% Beamer configuration
%

\usefonttheme{default}
\usefonttheme{serif}
\usefonttheme{structurebold}

\setbeamersize{text margin left=6mm}
\setbeamersize{text margin right=6mm}
\setbeamertemplate{headline}{}
\setbeamertemplate{navigation symbols}{}
\setbeamertemplate{footline}[page number]{}
\setbeamertemplate{itemize items}[circle]
\setbeamertemplate{theorems}[numbered]

\leftskip=0pt
\rightskip=0pt


\newcommand\Clang{\textsf{C}\xspace}
\newcommand{\Ada}{\textsf{Ada}\xspace}
\newcommand{\Java}{\textsf{Java}\xspace}
\newcommand{\OCaml}{\textsf{OCaml}\xspace}
\newcommand{\Maude}{\textsf{Maude}\xspace}
\newcommand{\Csharp}{\textsf{C}\raise 0.7mm \hbox{\(\scriptscriptstyle \textsf{\#}\)}\xspace}
\newcommand{\NaN}{\textsf{NaN}\xspace}
\newcommand{\true}{\textsf{true}\xspace}
\newcommand{\false}{\textsf{false}\xspace}
\newcommand\fun[1]{\textsf{#1}}
\newcommand\ufun[1]{\underline{\fun{#1}}}
\newcommand\el{[\,]}
\newcommand\cons[2]{{#1}\!::\!{#2}}
\newcommand\ttcons[2]{{#1}\,$::$\,{#2}}
\newcommand\FV[1]{\mathcal{F}({#1})}
\newcommand\smashedrightarrow[1]{\xrightarrow{\smash{#1}}}

\newcommand\measure[1]{\mathcal{M}\llbracket{#1}\rrbracket}

% Induction
%
\newcommand\predName[1]{\textsf{#1}}
\newcommand\pred[2]{\predName{#1}({#2})}
\newcommand\IH[1]{#1}

\title{A Theory of Programming}
\author{Christian Rinderknecht}
\institute{Eszterházy Károly Katolikus Egyetem}
\date{September 2026}

\begin{document}

\frame{\maketitle}

% ------------------------------------------------------------------------
%
\begin{frame}{Introduction}

  This lecture proposes a stepwise introduction to a foundation of
  programming and the analyses to write reliable and efficient
  software.

  \bigskip

  There are many different frameworks. We choose here a formal logic
  where relations and functions are the main building blocks. This
  entails that the meaning of \textbf{functional programs} can often
  be made very precise. (We will cover imperative programs later.)

  \bigskip

  It is a good framework to get a sense of \textbf{static analyses},
  and, as such, it stands a prominent place in the field of formal
  methods and \textbf{compiler construction}.

\end{frame}

% ------------------------------------------------------------------------
%
\begin{frame}{String rewriting}
  Let us consider a string of white and black beads, like \(\circ
  \bullet \bullet \bullet \circ \circ \bullet\), and a one-player game
  whose unique move consists in replacing two adjacent beads with only
  one according to the rules
  \begin{equation*}
    \bullet \; \circ   \xrightarrow{\alpha} \bullet\qquad\qquad
    \circ   \; \bullet \xrightarrow{\beta} \bullet\qquad\qquad
    \bullet \; \bullet \xrightarrow{\gamma} \circ
  \end{equation*}

  The rules~\(\alpha\), \(\beta\)~and~\(\gamma\) make up a simple
  \textbf{string\hyp{}rewriting system}.

  \bigskip

  Rules \(\alpha\)~and~\(\beta\) can be conceived as:\\ \emph{`A black
  bead absorbs the white bead next to it.'}

  \bigskip

  The goal of this game is to end up with as few beads as possible, so
  our example may lead to the \textbf{rewrites}
  \begin{equation*}
    \circ \bullet \bullet \, \fbox{\(\bullet \, \circ\)} \circ \bullet
    \xrightarrow{\alpha} \circ \bullet \bullet \, \fbox{\(\bullet \,
      \circ\)} \, \bullet \xrightarrow{\alpha} \fbox{\(\circ \,
      \bullet\)} \bullet \bullet \bullet \xrightarrow{\beta} \bullet
    \bullet \fbox{\(\bullet \, \bullet\)} \xrightarrow{\gamma} \bullet \,
    \fbox{\(\bullet \, \circ\)} \xrightarrow{\alpha} \fbox{\(\bullet
      \, \bullet\)} \xrightarrow{\gamma} \circ
  \end{equation*}

  \medskip

  \noindent where the part of the string to be rewritten next is
  framed.

\end{frame}

% ------------------------------------------------------------------------
%
\begin{frame}{Normal forms and termination}

  \label{termination1}

  Other compositions of the rules lead to the same result~\(\circ\) as
  well. Some others bring all\hyp{}white or all\hyp{}black strings,
  like \(\circ \, \circ\) and~\(\bullet\).

  \bigskip

  Strings that can not be further rewritten, or \textbf{reduced}, are
  called \textbf{normal forms}. We may wonder whether all strings have
  a normal form, and, if so, if it is unique and, furthermore, if it
  is either all\hyp{}white or all\hyp{}black.

  \bigskip

  First, let us note that the system is \textbf{terminating}, that is,
  there is no infinite chain of rewrites, because the number of beads
  strictly decreases in all the rules.

  \bigskip

  This is not a necessary condition in general, for instance, \(\circ
  \; \bullet \xrightarrow{\beta} \bullet \circ \circ\) would preserve
  termination because the composition \(\beta\alpha\alpha\) would be
  equivalent to the original rule~\(\beta\).

  \bigskip

  In particular, this means that any string has a normal form.

\end{frame}

% ------------------------------------------------------------------------
%
\begin{frame}{Invariants}

  Second, notice how the parity of the number of black beads is
  \textbf{invariant} through each rule and how there is no rewrite
  rule for two adjacent white beads.

  \bigskip

  Therefore, if there are \(2p\) initial black beads, then composing
  rules \(\alpha\)~and~\(\beta\) lead to an all\hyp{}black string,
  which can be reduced by applying rule~\(\gamma\) to contiguous pairs
  of beads into an all\hyp{}white string made of \(p\)~beads.

  \bigskip

  Otherwise, the same all\hyp{}black string can be reduced by applying
  alternatively \(\gamma\)~and~\(\beta\) on the left end or
  \(\gamma\)~and~\(\alpha\) on the right end, yielding~\(\circ\).

  \bigskip

  Similarly, if there is an initial odd number of black beads, we
  always end up with one black bead.

\end{frame}

% ------------------------------------------------------------------------
%
\begin{frame}{Confluence}

  Consequently, normal forms are not unique. A simple example is
  \begin{equation*}
    \circ \circ \xleftarrow{\gamma} \bullet \bullet \circ
    \xrightarrow{\alpha} \bullet \bullet \xrightarrow{\gamma}
    \circ
  \end{equation*}

  \medskip

  Systems where normal forms are unique are \textbf{confluent}.

  \bigskip

  If we think of a rewrite step as an elementary computation, it
  becomes clear that confluence is a desirable property for carrying
  out computations, as we would expect, for instance, an arithmetic
  expression to have the same value no matter the order of the small
  steps.

  \bigskip

  A systemic order of rule application is called a \textbf{reduction
    strategy}.

\end{frame}

% ------------------------------------------------------------------------
%
\begin{frame}{Completing a system}

  We can obtain confluence by adding the rule~\(\delta\):
  \begin{equation*}
    \bullet \; \circ   \xrightarrow{\alpha} \bullet\qquad\qquad
    \circ   \; \bullet \xrightarrow{\beta} \bullet\qquad\qquad
    \bullet \; \bullet \xrightarrow{\gamma}
    \circ\qquad\qquad
    \circ\; \circ \xrightarrow{\delta} \circ
  \end{equation*}
  The result of the game is now always one bead whose colour depends
  on the original parity of the black beads as before.

  \bigskip

  To see why, let us consider first that two non\hyp{}overlapping
  parts of a string can be rewritten in parallel, so they are not a
  concern.

  \bigskip

  The interesting cases occur when two applications of rules (maybe of
  the same rule) lead to different strings because they do
  overlap. For instance,
  \begin{equation*}
    \circ \, \circ \xleftarrow{\gamma} \bullet \bullet
    \circ \xrightarrow{\alpha} \bullet \, \bullet
  \end{equation*}

  \medskip

  The important point is that \(\circ \, \circ\) and \(\bullet \,
  \bullet\) can be rewritten into~\(\circ\) at the next step by
  \(\delta\)~and~\(\gamma\), respectively.

\end{frame}

% ------------------------------------------------------------------------
%
\begin{frame}{Critical pairs}

  In general, what matters is that all pairs of strings resulting from
  the application of overlapping rules, called \textbf{critical
    pairs}, can be rewritten to the same string: they are
  \textbf{joinable}.

  \bigskip

  In our example, all interactions occur on substrings made of three
  beads, so we must examine eight cases:
  \begin{center}
    \includegraphics[]{bullets}
  \end{center}

  In all the cases, the divergences are joinable in one step at most.

\end{frame}

% ------------------------------------------------------------------------
%
\begin{frame}{Local confluence}

  If all critical pairs are joinable after a divergence in one ore
  more rewrites, the system is \textbf{locally confluent}.

  \bigskip

  Moreover, if the system is also terminating, then \textbf{every
    string has exactly one normal form}.

  \bigskip

  Sometimes, as seen previously, a non\hyp{}confluent system can be
  completed to become confluent, via a \textbf{completion procedure},
  for example, in the case of terminating systems, the
  \textbf{Knuth\hyp{}Bendix semi\hyp{}algorithm}.

\end{frame}

% ------------------------------------------------------------------------
%
\begin{frame}{Ground and linear rewrite systems}

  The system we defined is \textbf{ground}, that is, it involves no
  variables.

  \bigskip

  These allow a finite system to denote an infinite number of ground
  rules if variables range over infinite sets, or simply, they reduce
  the size of rewrite systems.

  \bigskip

  For instance, our continued example is equivalent to
  \begin{equation*}
    \bullet \; \circ \xrightarrow{\alpha} \bullet
    \qquad\qquad
    \circ \; x \xrightarrow{\beta+\delta} x
    \qquad\qquad
    \bullet \; \bullet \xrightarrow{\gamma} \circ
  \end{equation*}

  \medskip

  If we accept multiple occurrences of a variable on the
  left\hyp{}hand side of a rule (that is, not left\hyp{}linear), we
  can further decrease the size of the system:
  \begin{equation*}
    x  \; x \xrightarrow{\gamma+\delta} \circ\qquad\qquad
    x  \; y \xrightarrow{\alpha+\beta} \bullet
  \end{equation*}

\end{frame}

% ------------------------------------------------------------------------
%
\begin{frame}{Ordered rewrite systems}

  Notice that there is now an implicit order over the rules: the rule
  \(\gamma+\delta\) must be examined first for a \textbf{match} with a
  part of the current string, because it is included in the second
  (set \(x=y\) in \(\alpha+\beta\)).

  \bigskip

  Non\hyp{}linear systems are equational theories due the implicit
  equality on substructures, which is complex: we will limit ourselves
  to \textbf{linear systems}.

  \bigskip

  In general, the order in which the rules are written on the page is
  their \textbf{logical order}.

  \bigskip

  Moreover, let us remark that the system does not specify that
  \(x\)~must either be~\((\circ)\) or~\((\bullet)\): this should be
  stated nevertheless somewhere else.

  \bigskip

  Generally speaking, this means that the \textbf{type} of the
  variables has to be defined elsewhere or inferred from the variable
  uses.

\end{frame}

% ------------------------------------------------------------------------
%
\begin{frame}{A simple calculator}

  Perhaps, after playing with beads, the simplest thing is an integer
  calculator, that is, integers and the four basic arithmetic
  operators.

  \bigskip

  An \textbf{expression} may be built using these operations, integers
  and parentheses. For example \((2+3)\times(4+5)\) is an expression.

  \bigskip

  A \textbf{value}  is an integer, therefore  it is a special  case of
  expression.

  \bigskip

  What a calculator does is to match an input expression with a value,
  the result. For example, the above expression evaluates in~\(45\).

  \bigskip

  This process is called \textbf{evaluation}.

\end{frame}

% ------------------------------------------------------------------------
%
\begin{frame}{Evaluation as rewriting expressions / Reductions}

  The evaluation of arithmetic expressions can be defined as a series
  of transformations leading to a value (the result). These
  transformations are called \textbf{rewrites} and a series of
  rewrites is a \textbf{reduction} or, more generally, a
  \textbf{computation}.

  \bigskip

  As we learnt in school long time ago, the way to perform these
  rewrites is to select a part of the current expression consisting of
  an operator and its arguments \emph{as values}, then to replace this
  sub-expression by the operation result.

  \bigskip

  In the following reduction, let us underline the sub-expression to
  be rewritten and print in bold the result of each step:
  \begin{equation*}
    (2+3) \times (\underline{4+5}) \rightarrow (\underline{2+3})
    \times \textbf{9} \rightarrow \underline{\textbf{5} \times 9}
    \rightarrow \textbf{45}
  \end{equation*}

\end{frame}

% ------------------------------------------------------------------------
%
\begin{frame}{Strategy of evaluation and rewrite systems}

  Note that there are different ways to choose the sub-expression. For
  example:
  \begin{equation*}
    (\underline{2+3}) \times (4+5) \rightarrow \textbf{5} \times
    (\underline{4+5}) \rightarrow \underline{5 \times \textbf{9}}
    \rightarrow \textbf{45}
  \end{equation*}
  It is usual to write, as a summary: \((2+3)\times(4+5)
  \xrightarrow{*} 45\) or \((2+3)\times(4+5)=45\).

  \bigskip

  So, there are two concepts involved here:
  \begin{enumerate}

    \item the selection of a sub-expression made of an operator and
      its arguments as values,

    \item the rewriting of such sub-expression into a value.

  \end{enumerate}
  They are defined by a \textbf{rewrite system}.

\end{frame}

% ------------------------------------------------------------------------
%
\begin{frame}{Rewrite rules}

  A rewrite system is a set of \textbf{rewrite rules}.

  \bigskip

  Rewriting rules specify how to perform a rewrite.

  \bigskip

  For example, let us assume that we have an infinite set of rewriting
  rules, that is. an infinite rewrite system, like:
  \begin{align*}
    0 + 0 & \rightarrow 0\\
    0 + 1 & \rightarrow 1\\
    \dots & \rightarrow \dots\\
    124 - 57 & \rightarrow 67\\
    \dots & \rightarrow \dots
  \end{align*}

\end{frame}

% ------------------------------------------------------------------------
%
\begin{frame}{Conditional rewrite rules and meta-variables}

  This rewrite system is not enough for the kind of expressions we
  want to evaluate, where it is possible to combine operations, like
  \((2+3) \times (4+5)\).

  \bigskip

  So let us extend our system with rules that specify the
  sub-expressions we can rewrite.

  \bigskip

  Consider the multiplication:
  \begin{equation*}
    e_1 \times e_2 \rightarrow e'_1 \times e_2 \quad \text{if} \; e_1
    \rightarrow e'_1
  \end{equation*}
  where \(e_1\), \(e'_1\) and \(e_2\) denote any expression and are
  called \textbf{meta\hyp{}variables}.

\end{frame}

% ------------------------------------------------------------------------
%
\begin{frame}{Instance of a rewrite rule}

  An \textbf{instance} of this rule is obtained when the
  meta-variables are replaced by expressions.

  \bigskip

  For example, \( (2+3) \times (4+5) \rightarrow 5 \times (4+5) \) is
  an instance since \(2 + 3 \rightarrow 5\). To understand why,
  replace \(e_1\) by \(2+3\), \(e'_1\) by \(5\) and \(e_2\) by
  \(4+5\).

  \bigskip

  Let us rewrite \(((2 + 3) \times 6) \times (4+5)\). First, we see
  that the previous rule has the required shape, i.e., the left side
  of the arrow is \(e_1 \times e_2\) and if \(e_1 = (2+3) \times 6\)
  and \(e_2 = 4+5\) then \(e_1 \times e_2\) is exactly the expression
  we wish to rewrite. So, we replace \(e_1\) and \(e_2\) by their
  associated expressions in the rule and we get an instance of it:
  \begin{equation*}
    ((2 + 3) \times 6) \times (4+5) \rightarrow e'_1 \times (4+5)
    \quad \text{if} \; (2+3) \times 6 \rightarrow e'_1
  \end{equation*}
  Therefore, we now know that we have to rewrite \((2 + 3) \times 6\).

\end{frame}

% ------------------------------------------------------------------------
%
\begin{frame}{Instance of a rewrite rule}

  We need another instance of the same rule. This time, let \(e_1 =
  2+3\) and \(e_2 = 6\). We get
  \begin{equation*}
    (2+3) \times 6 \rightarrow e'_1 \times 6 \quad \text{if} \; 2 + 3
    \rightarrow e'_1
  \end{equation*}
  It is important to understand that this latter \(e'_1\) is not the
  same as the former: they belong to two different instances of the
  same rule. The latter \(e'_1\) is actually easy to find, since we
  have all the rules for basic cases as \(2 + 3 \rightarrow
  5\). Hence, \(e'_1 = 5\) and the second instance is now complete:
  \begin{equation*}
    (2+3) \times 6 \rightarrow 5 \times 6 \quad \text{if} \; 2 + 3
    \rightarrow 5
  \end{equation*}
  Going backwards, this implies that the former \(e'_1\) is \(5 \times
  6\). So, the first instance is
  \begin{equation*}
    ((2+3) \times 6) \times (4+5) \rightarrow (5 \times 6) \times
    (4+5) \quad \text{if} \; 2 + 3 \rightarrow 5
  \end{equation*}

\end{frame}

% ------------------------------------------------------------------------
%
\begin{frame}{Instance of a rewrite rule}

  Our rule for multiplication, \(e_1 \times e_2 \rightarrow e'_1
  \times e_2\), actually depends on another, \(e_1 \rightarrow
  e'_1\). It is usual to write them vertically:
  \begin{mathpar}
    \inferrule*[right=\(\TirName{Mult}_1\)]
               {e_1 \rightarrow e'_1}
               {e_1 \times e_2 \rightarrow e'_1 \times e_2}
  \end{mathpar}
  This is called an \textbf{inference rule}. Note that it has a name,
  \(\TirName{Mult}_1\). The rewrite rule on the top, acting as a
  condition for the rule on the bottom, is called a
  \textbf{premise}. The rule at the bottom is called a
  \textbf{conclusion}.

  \bigskip

  A \textbf{substitution} is a mapping from meta-variables to
  expressions which applies to an inference rule or rewrite rule. For
  example, let \(\sigma\) be the following substitution: \(\{e_1
  \mapsto (2+3) \times 6; e_2 \mapsto 4+5\}\). This notation is
  equivalent to \(\sigma(e_1) = (2+3) \times 6\) and \(\sigma(e_2) =
  4+5\).

\end{frame}

% ------------------------------------------------------------------------
%
\begin{frame}{Instance of a rewrite rule}

  Now we can make a clear definition of the notion of `instance of a
  rule':

  \medskip

  \begin{quote}
    \emph{An instance of a rule is a rule in which some meta-variables
    \(e_i\) are replaced by expressions \(\sigma(e_i)\).}
  \end{quote}
  For example, here is the instance of \(\TirName{Mult}_1\) using
  previously defined substitution \(\sigma\), and noted
  \(\sigma(\TirName{Mult}_1)\):
  \begin{mathpar}
    \inferrule*[right=\(\sigma(\TirName{Mult}_1)\)]
               {(2 + 3) \times 6 \rightarrow e'_1}
               {((2+3) \times 6) \times (4+5) \rightarrow e'_1 \times (4+5)}
  \end{mathpar}

\end{frame}

% ------------------------------------------------------------------------
%
\begin{frame}{Instance of a rewrite rule}

  Now how do we find (a mapping for) \(e'_1\)?

  \bigskip

  We define another substitution \(\sigma'=\{e_1 \mapsto 2 + 3; e_2
  \mapsto 6\}\) and apply it:
  \begin{mathpar}
    \inferrule*[right=\(\sigma'(\TirName{Mult}_1)\)]
        {2 + 3 \rightarrow e'_1}
        {(2 + 3) \times 6 \rightarrow e'_1 \times 6}
  \end{mathpar}
  Note that the meta-variable \(e'_1\) in
  \(\sigma'(\TirName{Mult}_1)\) is not the same as in
  \(\sigma(\TirName{Mult}_1)\): they should have a different mapping.

  \bigskip

  We have a basic rewrite rule \(2 + 3 \rightarrow 5\), thus we extend
  \(\sigma'\) with \(e'_1 \mapsto 5\) and get
  \begin{mathpar}
    \inferrule*[right=\(\sigma'(\TirName{Mult}_1)\)]
        {2 + 3 \rightarrow 5}
        {(2 + 3) \times 6 \rightarrow 5 \times 6}
  \end{mathpar}

\end{frame}

% ------------------------------------------------------------------------
%
\begin{frame}{Instance of a rewrite rule}

  Therefore we can extend \(\sigma\) with \(e'_1 \mapsto 5 \times 6\)
  and we get
  \begin{mathpar}
    \inferrule*[right=\(\sigma(\TirName{Mult}_1)\)]
       {(2 + 3) \times 6 \rightarrow 5 \times 6}
       {((2+3) \times 6) \times (4+5) \rightarrow (5 \times 6) \times (4+5)}
  \end{mathpar}
  We can summarise this rewrite using the two instances of
  \(\TirName{Mult}_1\) as
  \begin{mathpar}
    \inferrule*[right=\(\sigma(\TirName{Mult}_1)\)] {
      \inferrule*[right=\(\;\sigma'(\TirName{Mult}_1)\)]
                 {2 + 3 \rightarrow 5}
                 {(2 + 3) \times 6 \rightarrow 5 \times 6}
    }
    {((2+3) \times 6) \times (4+5) \rightarrow (5 \times 6) \times (4+5)}
  \end{mathpar}

\end{frame}

% ------------------------------------------------------------------------
%
\begin{frame}{Instance of a rewrite rule}

  But what about the rewrite of an expression like \(7 \times (4+5)\)?

  \bigskip

  Let us define another substitution on \(\TirName{Mult}_1\):
  \(\sigma' = \{e_1 \mapsto 7; e_2 \mapsto 4+5\}\).

  \bigskip

  The corresponding instance is
  \begin{mathpar}
    \inferrule*[right=\(\sigma'(\TirName{Mult}_1)\)]
        {7 \rightarrow e'_1}
        {7 \times (4 + 5) \rightarrow e'_1 \times (4 +5)}
  \end{mathpar}
  But there is no rule of shape `\(v \rightarrow \dots\)', where \(v\)
  is a \emph{value} (because a value is, by definition, an expression
  completely reduced).

\end{frame}

% ------------------------------------------------------------------------
%
\begin{frame}{Completing the multiplication rule}

  This means that we need another rule for the evaluation of the
  \emph{right} argument of \((\times)\):
  \begin{mathpar}
    \inferrule*[right=\(\TirName{Mult}_2\)]
        {e_2 \rightarrow e'_2}
        {e_1 \times e_2 \rightarrow e_1 \times e'_2}
  \end{mathpar}
  This inference rule says that we can rewrite a product by rewriting
  its right argument.

  \bigskip

  For example, here is the instance \(\sigma'(\TirName{Mult}_2)\):
  \begin{mathpar}
    \inferrule*[right=\(\sigma'(\TirName{Mult}_2)\)]
        {4 + 5 \rightarrow 9}
        {7 \times (4+5) \rightarrow 7 \times 9}
  \end{mathpar}

\end{frame}

% ------------------------------------------------------------------------
%
\begin{frame}{Strategy}

  But now, something interesting happens.

  \bigskip

  Consider again \((2 + 3) \times (4+5)\): both \(\TirName{Mult}_1\)
  and \(\TirName{Mult}_2\) can be instantiated by substitution
  \(\sigma'' = \{e_1 \mapsto 2+3; e_2 \mapsto 4+5\}\):
  \begin{mathpar}
    \inferrule*[right=\(\sigma''(\TirName{Mult}_1)\)]
        {2 + 3 \rightarrow 5}
        {(2+3) \times (4+5) \rightarrow 5 \times (4+5)}
    \and
    \inferrule*[right=\(\sigma''(\TirName{Mult}_2)\)]
        {4 +5 \rightarrow 9}
        {(2+3) \times (4+5) \rightarrow (2+3) \times 9}
  \end{mathpar}

\end{frame}

% ------------------------------------------------------------------------
%
\begin{frame}{Strategy}

  This means that \emph{our system does not specify the order of
  evaluation of the arguments to \((\times)\).}

  \bigskip

  In other words, our rewrite system does not select only one
  sub\hyp{}expression to rewrite at each step.

  \bigskip

  We need a \textbf{strategy} to complete it.

\end{frame}

% ------------------------------------------------------------------------
%
\begin{frame}{Finiteness and soundness of reductions}

  At this point, it is clear that these reductions satisfy the
  following properties:
  \begin{itemize}

    \item \textbf{Finiteness.} All the reductions are finite, i.e.,
      all the computations terminate.

    \item \textbf{Soundness.} All the reductions of an expression end
      with the same value, if any. (Think confluence.)

  \end{itemize}

  About the second point, note indeed that \emph{some expressions have
  no value}, as
  \begin{equation*}
    (\underline{2+3})/(4-4) \rightarrow \textbf{5}/(\underline{4-4})
    \rightarrow 5/\textbf{0} \nrightarrow
  \end{equation*}
  The reduction is finite but does not end with a value: this
  situation corresponds to an \textbf{error}, usually called
  \textbf{run\hyp{}time error} in programming languages. Expression
  \(5/0\) is \textbf{stuck}.

\end{frame}

% ------------------------------------------------------------------------
%
\begin{frame}{Another model of run-time errors}

  We can modify the calculator so that, in case of error (here, the
  only one is the division by zero), the result of a rewrite is a
  special value, called \NaN (\emph{Not a Number}). Then we add the
  following rules:
  \begin{mathpar}
    \inferrule*[right=\(\quad\TirName{DivZero}\)]
       {}
       {e/0 \rightarrow \NaN}
    \and
    \inferrule*[right=\(\quad\TirName{Div-Err}_1\)]
       {}
       {\NaN/e \rightarrow \NaN}
    \and
    \inferrule*[right=\(\quad\TirName{Div-Err}_2\)]
       {}
       {e/\NaN \rightarrow \NaN}
    \and
    \inferrule*[right=\(\quad\TirName{Mult-Err}_1\)]
      {}
      {\NaN \times e \rightarrow \NaN}
    \and
    \inferrule*[right=\(\quad\TirName{Mult-Err}_2\)]
       {}
       {e \times \NaN \rightarrow \NaN}
    \and
    \inferrule*[right=\(\quad\TirName{Add-Err}_1\)]
       {}
       {\NaN + e \rightarrow \NaN}
    \and
    \inferrule*[right=\(\quad\TirName{Add-Err}_2\)]
       {}
       {e + \NaN \rightarrow \NaN}
    \and
    \inferrule*[right=\(\quad\TirName{Sub-Err}_1\)]
       {}
       {\NaN - e \rightarrow \NaN}
    \and
    \inferrule*[right=\(\quad\TirName{Sub-Err}_2\)]
       {}
       {e - \NaN \rightarrow \NaN}
\end{mathpar}

\end{frame}

% ------------------------------------------------------------------------
%
\begin{frame}{Another model of run-time errors}

  The rationale of this system is that once an error occurs, no other
  computations (leading to integer values) are required and the
  reduction can end as soon as possible.

  \bigskip

  In this case:
  \begin{equation*}
    (2+3)/(4-4) \rightarrow 5/(4-4) \rightarrow 5/0 \rightarrow \NaN
  \end{equation*}
  or, even the shortest possible:
  \begin{equation*}
    (2+3)/(4-4) \rightarrow (2+3)/0 \rightarrow \NaN
  \end{equation*}

\end{frame}

% ------------------------------------------------------------------------
%
\begin{frame}{Rephrasing soundness}

  The interesting phenomenon here is that, now, any reduction ends in
  a value (since errors are values). In the previous rewrite system,
  any reduction ends either in a value or a \textbf{stuck expression},
  i.e., an expression that cannot be rewritten further.

  \bigskip

  With the new system, that explicitly specifies the reduction of the
  \NaN error, the soundness property can be rephrased as

  \medskip

  \begin{quote}
    \emph{All the reductions of an expression end with the same value.}
  \end{quote}
  Obviously, the drawback of the new system is that it requires a lot
  of additional rules, and, as the system is augmented with new rules,
  new rules for the errors have to be added... or are forgotten by
  mistake.

\end{frame}

% ------------------------------------------------------------------------
%
\begin{frame}{Specifying the strategy}

  It is possible to force an order of evaluation on the arguments of
  \((\times)\). Assume we want to reduce the left argument first.

  \bigskip

  The idea is to change \(\TirName{Mult}_2\) so that instead
  of~\(e_1\), which stands for any expression, we specify a
  value~\(v\):
  \begin{mathpar}
    \inferrule*[right=\(\TirName{New-Mult}_2\)]
       {e_2 \rightarrow e'_2}
       {v \times e_2 \rightarrow v \times e'_2}
  \end{mathpar}
  This way, it is not possible anymore to instantiate this rule with
  expression \((2+3) \times (4+5)\) because \(2+3\) is not a value.

\end{frame}

% ------------------------------------------------------------------------
%
\begin{frame}{Specifying the strategy}

  The only possible reduction is:
  \begin{mathpar}
    \inferrule*[right=\(\sigma_1(\TirName{Mult}_1)\)]
       {2 + 3 \rightarrow 5}
       {(2+3) \times (4+5) \rightarrow 5 \times (4+5)}
    \and
    \inferrule*[right=\(\sigma_2(\TirName{New-Mult}_2)\)]
       {4 + 5 \rightarrow 9}
       {5 \times (4+5) \rightarrow 5 \times 9}
  \end{mathpar}
  where \(\sigma_1=\{e_1 \mapsto 2+3; e_2 \mapsto 4+5\}\) and
  \(\sigma_2=\{v \mapsto 5; e_2 \mapsto 4+5\}\).

\end{frame}

% ------------------------------------------------------------------------
%
\begin{frame}{Determinism}

  We have specified the strategy \textbf{in} the rewrite system.

  \bigskip

  Also, with \(\TirName{New-Mult}_2\), given an expression, there is
  only one rule that can be applied (or none): this is sometimes
  called \textbf{determinism}.

  \bigskip

  \textbf{Determinism.} If \(e_1 \rightarrow e'_1\) and \(e_2
  \rightarrow e'_2\) then \(e'_1 = e'_2\).

  \bigskip

  This is a stronger property than soundness, i.e., it implies
  soundness. Indeed, determinism implies that there is only one
  reduction from an expression to its value (if any).

\end{frame}

% ------------------------------------------------------------------------
%
\begin{frame}{Pragmatics of programming languages}

  In many programming languages, the order of evaluation of operator
  arguments is not specified, which usually means that there is a
  sentence in the official document which defines the language like:
  `The order of evaluation of the arguments of operators and functions
  is not specified.'

  \bigskip

  Formally, this means that the underlying rewrite system is
  \textbf{non-deterministic}.

  \bigskip

  One notable exception is \Java, which specifies that the arguments
  are always evaluated from left to right, i.e., following the order
  of writing.

  \bigskip

  Of course, if an order is specified, as in \Java, the compilers of
  these languages must implement it. Otherwise one is chosen
  arbitrarily and may change from one release to another.

\end{frame}

% ------------------------------------------------------------------------
%
\begin{frame}{Rewrite systems as models of languages}

  By studying rewrite systems, we can learn the basic concepts
  involved in programming languages, because rewrite systems are
  extremely simple.

  \bigskip

  There are some languages, like \Maude, whose programs actually are
  rewrite systems.

  \bigskip

  Some compilers transform the input program (called \textbf{source})
  and output a program, called \textbf{byte-code}, which can be
  reduced by a rewrite system.

\end{frame}

% ------------------------------------------------------------------------
%
\begin{frame}{Rewrite systems as models of languages}

  This is the case of \Java. This is why, in general, we need a
  \textbf{\Java virtual machine} (JVM) to run the program: this JVM
  performs the reductions according to a rewrite system.

  \bigskip

  Even if there is no byte-code, we always can interpret what the
  program does by studying a rewrite system which models it.

\end{frame}

% ------------------------------------------------------------------------
%
\begin{frame}{Sequentiality and parallelism}

  Some programming languages, like \Ada, allow the programmer to
  specify that different parts of its program can actually be executed
  separately: this is called, in general, \textbf{parallelism}. This
  parallelism is possible only if the underlying rewrite system is
  non-deterministic.

  \bigskip

  Indeed, if all the strategies lead to the same value (soundness
  property), it is possible to reduce different parts of an expression
  in parallel if the rewrite system allows it.

\end{frame}

% ------------------------------------------------------------------------
%
\begin{frame}{Sequentiality and parallelism}

  For example, the first argument of \((\times)\) can be computed
  \emph{at the same time} as its second argument. In the case of
  \((2+3) \times (4+5)\), this means that we can reduce in parallel
  \(2+3\), with \(\TirName{Mult}_1\), and \(4+5\), with
  \(\TirName{Mult}_2\), then \(5 \times 9\) is finally rewritten in
  \(45\) (after \textbf{synchronisation} of the two reductions).

\end{frame}

%------------------------------------------------------------------------
%
\begin{frame}[fragile]{Order of evaluation in C}

  In case of C, the order of evaluation arguments of arithmetic
  operators and functions is not specified, it actually depends on the
  compiler and its version. Therefore it would be dangerous to rely on
  the order of evaluation in C.

  \bigskip

  Why is the order left unspecified? Simply because of possible
  \textbf{side-effects}. Imagine the following situation:
{\small
\begin{verbatim}
int n = 0;
int f () { n++; return n; };
int g () { n = 2; return n; };
int main () { return f() + g() + n; }
\end{verbatim}
}
What is the result?

\end{frame}

% ------------------------------------------------------------------------
%
\begin{frame}{Boolean expressions}

  Let~\(\wedge\) represent the conjunction (`and'), \(\vee\) the
  disjunction (`or'), \(\neg\) the negation (`not'), and let the
  two values \true and \false.

  \bigskip

  The definition of the conjunction of values is very simple:
  \begin{mathpar}
    \inferrule
       {}
       {\true \wedge \true \rightarrow \true}
    \and
    \inferrule
       {}
       {\true \wedge \false \rightarrow \false}\\
    \inferrule
       {}
       {\false \wedge \true \rightarrow \false}
    \and
    \inferrule
       {}
       {\false \wedge \false \rightarrow \false}
\end{mathpar}

\end{frame}

% ------------------------------------------------------------------------
%
\begin{frame}{Boolean expressions made simpler}

  This system can even be made simpler, by means of meta-variables:
  \begin{mathpar}
    \inferrule*[right=\(\quad\TirName{True}\)]
       {}
       {\true \wedge \true \rightarrow \true}\\
    \inferrule*[right=\(\quad\TirName{False}_1\)]
       {}
       {\false \wedge v \rightarrow \false}\\
    \inferrule*[right=\(\quad\TirName{False}_2\)]
       {}
       {v \wedge \false \rightarrow \false}
  \end{mathpar}
  where \(v\) denotes any value, that is, either \true or \false.

\end{frame}

% ------------------------------------------------------------------------
%
\begin{frame}{Boolean expressions}

  So, if we define the rewrite of boolean expressions (not just the
  values \true and \false) as
  \begin{mathpar}
    \inferrule*[right=\(\TirName{\(\wedge_1\)}\)]
       {e_1 \rightarrow e'_1}
       {e_1 \wedge\, e_2 \rightarrow e'_1 \wedge\, e_2}
    \and
    \inferrule*[right=\(\TirName{\(\wedge_2\)}\)]
       {e_2 \rightarrow e'_2}
       {e_1 \wedge\, e_2 \rightarrow e_1 \wedge\, e'_2}
  \end{mathpar}
  then, for example,
  \begin{mathpar}
    \inferrule*[right=\(\sigma_1(\wedge_1)\)]
       {\false \wedge\; \true \rightarrow \false
        \quad \TirName{False}_1}
       {(\false \wedge\; \true) \wedge\, (\true \wedge\, \true)
         \rightarrow \false \wedge\, (\true \wedge\, \true)}
  \end{mathpar}
  where \(\sigma_1=\{e_1 \mapsto \false \wedge \true; e_2 \mapsto
  \true \wedge \true; e'_1 \mapsto \false\}\).

\end{frame}

% ------------------------------------------------------------------------
%
\begin{frame}{Boolean expressions}

  Then
  \begin{mathpar}
    \inferrule*[right=\(\sigma_2(\wedge_2)\)]
       {\true \wedge\, \true \rightarrow \true \quad \TirName{True}}
       {\false \wedge\, (\true \wedge\, \true) \rightarrow \false \wedge\,
  \true}
  \end{mathpar}
  where \(\sigma_2 = \{e_1 \mapsto \false; e_2 \mapsto \true \wedge
  \true; e'_1 \mapsto \true\}\).

  \bigskip

  and, finally, we can instantiate \(\TirName{False}_1\) with
  \(\sigma_3 = \{v \mapsto \true\}\):
  \begin{mathpar}
    \inferrule*[right=\(\quad\sigma_3(\TirName{False}_1)\)]
       {}
       {\false \wedge\, \true \rightarrow \false}
  \end{mathpar}
  In this example, we computed both arguments of \(\wedge\) to reach
  the value.

\end{frame}

% ------------------------------------------------------------------------
%
\begin{frame}{Boolean expressions}

  But the inspection of the basic rules shows that if one of the
  arguments of \(\wedge\) reduces to \false, then there is no need to
  reduce the other argument: the rewrite will always be
  \false. Therefore, it is more efficient to extend the rules on
  values to the expressions:
  \begin{mathpar}
    \inferrule*[right=\(\quad\TirName{False}_1\)]
       {}
       {\false \wedge e \rightarrow \false}
    \and
    \inferrule*[right=\(\quad\TirName{False}_2\)]
       {}
       {e \wedge \false \rightarrow \false}\\
    \inferrule*[right=\(\quad\TirName{True}_1\)]
       {}
       {\true \wedge e \rightarrow e}
    \and
    \inferrule*[right=\(\quad\TirName{True}_2\)]
       {}
       {e \wedge \true \rightarrow e}
  \end{mathpar}
  where \(e\) represents any boolean expression. Then, now
  \begin{mathpar}
    \inferrule*[right=\(\quad\sigma_3(\TirName{False}_1)\)]
       {}
       {\false \wedge\, (\true \wedge\, \true) \rightarrow \false}
  \end{mathpar}
  where \(\sigma_3 = \{e \mapsto \true \wedge \true\}\).

\end{frame}

% ------------------------------------------------------------------------
%
\begin{frame}[fragile]{Boolean expressions}

  Try to reduce
  \begin{equation*}
    (\true \wedge (\false \wedge \true)) \wedge (\true \wedge (\true
    \wedge \true))
  \end{equation*}
  Note that this rewrite system is not deterministic. If we enforce
  that the left argument of \(\wedge\) is completely reduced first,
  then we can check if it is \true or \false: in the latter case, we
  do not need to reduce the right argument.

  \begin{mathpar}
    \inferrule*[right=\(\quad\wedge_{\TirName{True}}\)]
       {}
       {\true \wedge e \rightarrow e}
    \and
    \inferrule*[right=\(\quad\wedge_{\TirName{False}}\)]
       {}
       {\false \wedge e \rightarrow \false}\\
    \inferrule*[right=\(\wedge\)]
       {e_1 \rightarrow e'_1}
       {e_1 \wedge\, e_2 \rightarrow e'_1 \wedge\, e_2}
  \end{mathpar}
  Many programming languages, as C, use this strategy. It allows the
  programmers to write handy code like
{\footnotesize
\begin{verbatim}
while (index < max && tab[index] > 0) ...
\end{verbatim}
}

\end{frame}

% ------------------------------------------------------------------------
%
\begin{frame}{Boolean expressions}

  Note that this strategy (left argument before the right one) is not
  optimal in general: the first argument may reduce to
  \true. \textbf{Finding an optimal strategy is not always possible.}
  For example, we could add another rule that shortens sometimes the
  reductions:
  \begin{mathpar}
    \inferrule*[right=\(\wedge_{\TirName{Refl}}\)]
       {}
       {e \wedge e \rightarrow e}
  \end{mathpar}
  This rule explicitly specifies the \textbf{reflexivity} of the
  conjunction.

  But the system becomes non-deterministic because there is now an
  overlap between rule \(\wedge_{\TirName{Refl}}\) and both
  \(\wedge_{\proc{True}}\) and \(\wedge_{\proc{False}}\) (consider for
  example \(\true \wedge \false\)).

\end{frame}

% ------------------------------------------------------------------------
%
\begin{frame}{Determinism versus non-determinism}

  Let us summarise the concepts of determinism and
  non\hyp{}determinism.
  \begin{itemize}

    \item Determinism means that, for any expression, there is only
      one rule to rewrite it. This implies that, every time a program
      is run with the same input, the output will be the same. This is
      a useful property, called soundness, in particular for
      debugging.

    \item Non-determinism means that, for any expression, there may be
      several rules to rewrite it. Moreover,
      \begin{itemize}

         \item if the rewrite system is sound, the output will always
           be the same for any run on the same input, even if the
           trace may be different each time, which makes the debugging
           more difficult, but enable parallel reductions.

         \item the rewrite system may be unsound.

      \end{itemize}
  \end{itemize}

\end{frame}

% ------------------------------------------------------------------------
%
\begin{frame}[fragile]{Determinism versus non-determinism}

  Some programming languages have features that may be both
  non\hyp{}deterministic and unsound.

  \bigskip

  Perhaps the most infamous case is in C, where variables do not need
  to be initialised at their declaration:
{\footnotesize
\begin{verbatim}
int a;
\end{verbatim}
}
  If the programmer forgets to give a value to the variable, any access
  to it is unspecified.

  \bigskip

  Practically, this means that whatever is located in the memory at
  the location of the variable will be used at each read access. Of
  course, this value may change from one run of the program to
  another, even on the same input.

\end{frame}

% ------------------------------------------------------------------------
%
\begin{frame}{Commutativity and its consequences}

  We may try another rewrite system for \(\wedge\):
  \begin{mathpar}
    \inferrule*[right=\(\quad\TirName{True}\)]
       {\true \wedge e \rightarrow e}
       {}
    \and
    \inferrule*[right=\(\quad\TirName{False}\)]
       {}
       {\false \wedge e \rightarrow e}\\
    \inferrule*[right=\(\wedge\)]
       {e_1 \rightarrow e'_1}
       {e_1 \wedge e_2 \rightarrow e'_1 \wedge e_2}
    \and
    \inferrule*[right=\(\TirName{Comm}\)]
       {}
       {e_1 \wedge e_2 \rightarrow e_2 \wedge e_1}
  \end{mathpar}
  While this system is in theory fine, there are two practical
  problems:
  \begin{enumerate}

    \item it is \textbf{not deterministic} due to \(\TirName{Comm}\)
      that overlaps with all the other rules;

    \item it may \textbf{not terminate} due to rule
      \(\TirName{Comm}\):
      \begin{align*}
        \true \wedge \false
        & \;\xrightarrow{\TirName{Comm}}\;
        \false \wedge \true \;\xrightarrow{\TirName{Comm}}\; \true
        \wedge \false\\
        & \;\xrightarrow{\TirName{Comm}}\; \false
        \wedge \true \rightarrow \cdots
      \end{align*}

  \end{enumerate}

\end{frame}

% ------------------------------------------------------------------------
%
\begin{frame}[fragile]{Non-termination}

  Now some reductions may not be finite, which models
  \textbf{non\hyp{}termination}.

  \bigskip

  This particular kind of non\hyp{}termination, due to commutativity,
  is the model in programming languages of so\hyp{}called
  `\textbf{infinite loops}' that do not allocate memory. Equivalently,
  it corresponds to a \emph{recursive call} like
{\footnotesize
\begin{verbatim}
bool and(bool e1, bool e2) { ... return and(e2, e1); ... }
\end{verbatim}
}
  Anyway, in the presence of infinite reductions, the soundness
  property must be weakened as follows:

  \bigskip

  \begin{quote}
    \textbf{Weak soundness.} \emph{All finite reductions of an
    expression end with the same value, if any.}
  \end{quote}

\end{frame}

% ------------------------------------------------------------------------
%
\begin{frame}{Non-termination}

  Let us define the negation:
  \begin{mathpar}
    \inferrule
       {}
       {\neg \true \rightarrow \false}
    \and
    \inferrule
       {}
       {\neg \false \rightarrow \true}
  \end{mathpar}
  It is also well known that the negation satisfies \(e = \neg\neg
  e\). So we could imagine another rewrite rule
  \begin{equation*}
    \neg\neg e \rightarrow e
  \end{equation*}
  This is fine, this rule simplifies double negations. But what if we
  had the converse rule?
  \begin{equation*}
    e \rightarrow \neg\neg e \qquad \text{where} \; e \; \text{is not a
  value.}
  \end{equation*}
  This allows infinite reductions: \(e \rightarrow \neg\neg e
  \rightarrow e \rightarrow \neg\neg e \rightarrow \dots\)

% excluded middle: e or ~e
% non-contradiction: ~(e and ~e)
% de Morgan applied to non-contradiction: ~(e and ~e) = ~e or ~~e

\end{frame}

% ------------------------------------------------------------------------
%
\begin{frame}{Summary}

  There are three kinds of reductions:
  \begin{enumerate}

    \item Finite and ending in a value
  \[e_1 \rightarrow e_2 \rightarrow \dots \rightarrow v\]

    \item Finite but ending in a stuck expression
  \[e_1 \rightarrow e_2 \rightarrow \dots \rightarrow e_n
    \nrightarrow\]

    \item Infinite
  \[e_1 \rightarrow e_2 \rightarrow \cdots\]

  \end{enumerate}

\end{frame}

% ------------------------------------------------------------------------
%
\begin{frame}{Conditional expressions}

  At this point, boolean expressions can be
  \begin{itemize}

    \item \Xtrue (a value)

    \item \Xfalse (a value)

    \item \(e_1 \wedge e_2\)

    \item \(e_1 \vee e_2\)

    \item \(\neg e_1\)

  \end{itemize}
  where \(e_1\) and \(e_2\) are boolean expressions. Let us extend our
  boolean expressions with a \textbf{conditional expression}:
  \begin{itemize}

     \item \Xif \(e_1\) \Xthen \(e_2\) \Xelse \(e_3\)

  \end{itemize}

\end{frame}

% ------------------------------------------------------------------------
%
\begin{frame}{Conditional expressions}

  It is easy to extend the rewrite system to cope with conditionals:
  \begin{mathpar}
    \inferrule*[right=\(\quad\TirName{If-True}\)]
       {}
       {\Xif \; \Xtrue \; \Xthen \; e_2 \; \Xelse \; e_3 \rightarrow e_2}
    \and
    \inferrule*[right=\(\quad\TirName{If-False}\)]
       {}
       {\Xif \; \Xfalse \; \Xthen \; e_2 \; \Xelse \; e_3 \rightarrow e_3}
    \and
    \inferrule*[right=\(\TirName{If}\)]
       {e_1 \rightarrow e'_1}
       {\Xif \; e_1 \; \Xthen \; e_2 \; \Xelse \; e_3
        \rightarrow \Xif \;
        e'_1 \; \Xthen \; e_2 \; \Xelse \; e_3}
  \end{mathpar}
  This means that the boolean expression \(e_1\) must be completely
  reduced \emph{before} the others. This is the usual way in
  programming languages, the rationale being to avoid unnecessary
  computations (evaluating \(e_2\) and~\(e_3\) is not necessary).

\end{frame}

% ------------------------------------------------------------------------
%
\begin{frame}{Term rewriting}

  More general are the \textbf{term\hyp{}rewriting systems}, where a
  \textbf{term} is a mathematical object built on tuples, integers and
  variables.

  \bigskip

   Let us consider the following totally ordered system:
   \begin{align*}
     (0,m) &\xrightarrow{1} m\\
     (n,m) &\xrightarrow{\smash 2} (n-1, n \times m)\\
     n     &\xrightarrow{\smash 3} (n,1)
   \end{align*}

   Arithmetic operators \((-)\)~and~\((\cdot)\) are defined out of the
   system and \(m\)~and~\(n\) are variables denoting natural numbers.

   \bigskip

   Would the rules not be ordered as they are actually laid out, the
   second rule would match any pair. Instead, it can be assumed that
   \(n \neq 0\) when matching with it.

   \bigskip

   The last rule is to be tried last, as it matches all terms.

\end{frame}

% ------------------------------------------------------------------------
%
\begin{frame}{Transitive closures}

  We can easily see that all the compositions of rewrites starting
  with a natural number~\(n\) end with the factorial of~\(n\):
  \smallskip
  \begin{equation*}
    n \xrightarrow{3} (n,1)
    \xrightarrow{2} \dots
    \xrightarrow{2} (0,n!)
    \xrightarrow{1}
    n!, \quad \text{for \(n \in \mathbb{N}\)}.
  \end{equation*}

  Let us note \((\xrightarrow{n})\) the composition of
  \((\rightarrow)\) repeated \(n-1\)~times:
  \begin{align*}
    (\xrightarrow{1})   &:= (\rightarrow)\\
    (\xrightarrow{\smash{n+1}}) & :=
    (\rightarrow) \circ (\xrightarrow{\smash{n}}),
    \quad\text{with \(n > 0\)}.
  \end{align*}

  The \textbf{transitive closure} of \((\rightarrow)\) is defined as
  \begin{equation*}
    (\twoheadrightarrow) := \bigcup_{i > 0}{(\xrightarrow{i})}.
  \end{equation*}

  The factorial function coincides with the transitive closure of
  \((\rightarrow)\): \(n \twoheadrightarrow n!\). Let
  \((\xrightarrow{*})\) be the reflexive\hyp{}transitive closure of
  \((\rightarrow)\), that is,
  \begin{equation*}
    (\xrightarrow{*}) := (=) \cup (\twoheadrightarrow).
  \end{equation*}

\end{frame}

% ------------------------------------------------------------------------
%
\begin{frame}{Functions}

  \begin{itemize}

    \item A confluent system defines a \textbf{partial function}.

    \item If the system is also terminating, we get a
      \textbf{function}.

    \item We can name functions: for example, \(\fun{c}(1,
      \fun{d}(n))\)~is a term constructed with \textbf{function names}
      \fun{c}~and~\fun{d}, as well as variable~\(n\).

    \item A tuple tagged with a function name, like \(\fun{f}(x,y)\),
      is called a \textbf{function call}.

    \item The components of the tuples are then called
      \textbf{arguments}, for example \(\fun{d}(n)\) is the second
      argument of the call \(\fun{c}(1, \fun{d}(n))\).

    \item It is possible for a function call to hold no arguments,
      like \(\fun{d}()\).

    \item We restrict the left\hyp{}hand sides of rules to be function
      calls.

  \end{itemize}

\end{frame}

% ------------------------------------------------------------------------
%
\begin{frame}{Trees}

  The topological understanding of a function call or a tuple is the
  finite \textbf{tree}. A tree is a hierarchical layout of
  information, for example:
  \begin{center}
    \includegraphics{tree_for_term}
  \end{center}

   The disks are called \textbf{nodes} and the segments which connect
   two nodes are called \textbf{edges}. The topmost node (with a
   diameter) is called the \textbf{root} and the bottommost ones
   (\(\bullet\)) are called the \textbf{leaves}. All nodes except the
   leaves are seen to downwardly connect to some other nodes, called
   \textbf{children}. Upwardly, each node but the root is connected to
   another node, called its \textbf{parent}.

\end{frame}

% ------------------------------------------------------------------------
%
\begin{frame}{Trees}

  \textbf{Trees are pervasive} in informatics, especially in
  functional programming, but not only. Structured documents, like
  books, articles, and web pages, are composed of chapters, sections,
  paragraphs, figures, appendices, indices, etc.

  \bigskip

  The occurrences of these components are mutually constrained; for
  instance, it is understood that a section is part of a chapter and
  that appendices are located at the end of a document.

  \bigskip

  This hierarchical layout is meant to facilitate reading, and it
  supports the search for specific items of information.

  \bigskip

  When considering computer systems, these data must be uniformly
  encoded by means of a formal language.

\end{frame}

% ------------------------------------------------------------------------
%
\begin{frame}{Terms as trees}

  \label{terms_as_trees}

  Trees can be used to depict terms as follows. A function call is a
  tree whose root is the function name and the children are the trees
  denoting the arguments.

  \bigskip

  A tuple can be considered as having an invisible function name
  represented by a node with a period (\texttt{.}) in the tree, in
  which case the components of the tuple are its children. For
  example,
  \begin{center}
    \includegraphics{tree}
  \end{center}
  is the tree corresponding to the tern
  \(\fun{f}((\fun{g}(0),(x,1)),(),\fun{g}(y))\).

  \bigskip

  Note that tuples, except the empty tuple, are not represented in the
  tree, since they encode the structure itself, which is already laid
  out. The number of arguments of a function is called \textbf{arity}.

\end{frame}

% ------------------------------------------------------------------------
%
\begin{frame}{Functional languages}

  \begin{itemize}

    \item For programming, we only consider confluent systems because
      they define partial functions.

    \item To have only one reduction chain per term to the normal
      form, we may also enforce that rules are ordered.

    \item We also enforce that normal forms must be \textbf{values},
      that is, terms without function calls or only with calls to
      functions not occurring on the left\hyp{}hand sides.

    \item In other words, irreducible calls of defined functions are
      not allowed in values.

    \item Systems with these constraints are called \textbf{functional
      rewrite systems}.

    \item A program here is a term, and evaluating it means to rewrite
      the term
      \begin{enumerate}
        \item to a value,
        \item or to end in error with a normal form that is not a value,
        \item or to never terminate.
      \end{enumerate}

  \end{itemize}

\end{frame}

% ------------------------------------------------------------------------
%
\begin{frame}{Call-by-value}

  If a function call contains at least one other function call, there
  is a choice as to which one to reduce first.

  \bigskip

  One reduction strategy, named \textbf{call-by-value}, consists in
  always reducing the arguments before the call itself. Unfortunately,
  call-by-value enables otherwise terminating rewrites to not
  terminate. For instance, let us consider
  \begin{equation*}
    \fun{f}(x) \xrightarrow{\alpha} 0\qquad
    \fun{g}() \xrightarrow{\beta} \fun{g}()
  \end{equation*}
  We have \(\fun{f}(\fun{g}()) \xrightarrow{\alpha} 0\) but
  \(\fun{f}(\fun{g}()) \xrightarrow{\beta} \fun{f}(\fun{g}())
  \xrightarrow{\beta} \dots\)

  \bigskip

  Despite the loss of expressiveness, we shall retain in this lecture
  the call by value, as it is the strategy of \OCaml and it is easy to
  follow.

  \bigskip

  In general, the order in which the arguments of a function call are
  evaluated is not specified.

\end{frame}

% ------------------------------------------------------------------------
%
\begin{frame}{Call-by-value}

  Another reason to choose call-by-value is that it allows us to
  restrict the shape of the left\hyp{}hand sides, called
  \textbf{patterns}, to one, outermost function call.

  \bigskip

  For instance, we can then disallow, for being useless, a
  rule like

  \begin{equation*}
    \fun{plus}(x,\fun{plus}(y,z)) \rightarrow
    \fun{plus}(\fun{plus}(x,y),z).
  \end{equation*}

  \noindent because, following call\hyp{}by\hyp{}value, the argument
  \(\fun{plus}(y,z)\) must be evaluated before the outermost call,
  therefore it cannot be part of any pattern.

\end{frame}

% ------------------------------------------------------------------------
%
\begin{frame}{Termination}

  \label{termination2}

  We do not require by construction that all functional systems
  terminate, although that is a desirable property, but we
  nevertheless speak of `functions' instead of `partial functions'
  (which would be technically correct).

  \bigskip

  If we enforce termination, we loose
  \textbf{Turing\hyp{}completeness}: we cannot program all computable
  functions.

  \bigskip

  The termination of rewrite systems and, in particular, functional
  programs, is an \textbf{undecidable problem} in general.

  \bigskip

  Nevertheless, there exists a set of well\hyp{}known \textbf{decision
    criteria}, which we will use when programming recursive functions
  in \OCaml.

\end{frame}

% ------------------------------------------------------------------------
%
\begin{frame}{Higher-order programs, recursion, $\lambda$-calculus}

  Most functional languages allow \textbf{higher\hyp{}order function}
  definitions, whereas standard term\hyp{}rewriting systems do
  not. For example:
  \begin{equation*}
    \fun{f}(g,0) \rightarrow 1
    \qquad
    \fun{f}(g,n) \rightarrow n \times g(g,n-1)
    \qquad
    \fun{fact}(n) \rightarrow \fun{f}(\fun{f},n)
  \end{equation*}
  where \(n \in \mathbb{N}\). Note that these two definitions are
  \textbf{not} recursive, yet the function \fun{fact} is the factorial
  function. A function definition is \textbf{recursive} if the name of
  the definition occurs on at least one right\hyp{}hand side.

  \bigskip

  An adequate theoretical framework to understand higher\hyp{}order
  functions is \textbf{\(\lambda\)-calculus}. In fact,
  \(\lambda\)-calculus features prominently in the semantics of
  programming languages, even not functional ones.

  \bigskip

  For didactic purposes, we began with rewrite systems because they
  offer implicit pattern matching (that is, rule selection).

\end{frame}

% ------------------------------------------------------------------------
%
\begin{frame}{Stacks in a functional setting}

  We show how to express linear data structures in a purely functional
  language and how to compute with them. The simplest of those
  structures is the \textbf{stack}.

  \bigskip

  A stack can be thought of as a finite series of items that can only
  be accessed sequentially from one end, called the \textbf{top},
  following the analogy with a stack of material objects.

\end{frame}

% ------------------------------------------------------------------------
%
\begin{frame}{Stacks in a functional setting}

  To define a stack in a functional language, one uses calls to
  undefined functions to model all the possible shapes a stack can
  take.

  \bigskip

  Here, a stack can be either empty or not.

  \bigskip

  Let \(\fun{nil}()\) be the irreducible call that denotes an empty
  stack (that is, the function `\fun{nil}' is undefined).

  \bigskip

  If not empty, the stack then must contain a topmost (first) item,
  and a \textbf{substack}, that is, the stack made of all the items
  but the first.

  \bigskip

  Let \(\fun{cons}(x,s)\) denote the non\hyp{}empty stack with the
  item~\(x\) on top and the \textbf{immediate substack}~\(s\).

\end{frame}

% ------------------------------------------------------------------------
%
\begin{frame}{Inductive definition of stacks}

  \begin{itemize}

    \item The previous definition is intuitive, but not formal yet.

    \item Data structures are usually defined \textbf{inductively}.

    \item Let~\(\mathcal{T}\) be the set of all possible terms and
      \(\mathcal{S} \subset \mathcal{T}\) be the set of all stacks.

    \item Formally, \(\mathcal{S}\)~can by defined by
      \textbf{induction} as the smallest set such that
      \begin{enumerate}

        \item \(\fun{nil}() \in \mathcal{S}\);

        \item if \(x \in \mathcal{T}\) and \(s \in \mathcal{S}\), then
          \(\fun{cons}(x,s) \in \mathcal{S}\).
      \end{enumerate}

    \item We take the smallest set because we do not want
      in~\(\mathcal{S}\) terms that were not constructed by the
      inductive rules (this is called the \textbf{smallest fixed
        point} of the inductive relation).

  \end{itemize}

\end{frame}

% ------------------------------------------------------------------------
%
\begin{frame}{Appending a stack}

  \begin{itemize}

    \item For all \(s, t \in \mathcal{S}\), let us consider the
      functional program
      \begin{align*}
        \fun{cat}(\fun{nil}(),t) &\xrightarrow{\alpha} t\\
        \fun{cat}(\fun{cons}(x,s),t) &\xrightarrow{\smash\beta}
        \fun{cons}(x,\fun{cat}(s,t))
      \end{align*}

    \item This is a functional program suitable for
      call\hyp{}by\hyp{}value reduction strategy, because the function
      calls in the patterns (that is, left\hyp{}hand sides) are
      irreducible.

    \item The right\hyp{}hand sides are stacks, so a call to
      `\fun{cat}' is a stack.

    \item Rule~\(\alpha\) says that, if the first stack is empty, the
      value is the second.

    \item In rule~\(\beta\), the first stack is not empty, and its top
      item is the top of the final stack~(\(x\)), the immediate
      substack~(\(s\)) being used in a call to `\fun{cat}' with the
      second stack~(\(t\)) unchanged.

  \end{itemize}

\end{frame}

% ------------------------------------------------------------------------
%
\begin{frame}{Appending a stack}

  This function appends the second stack at the bottom of a
  \textbf{copy} of the first. (The function name `\fun{cat}' stands
  for `catenation', and is slightly better than `\fun{app}').

  \bigskip

  We could also say that this function appends a copy of the first
  stack on top of the other. But this is not the same as pushing the
  first stack on top of the other, like
  \begin{equation*}
    \fun{cat}(s,t) \rightarrow \fun{cons}(s,t) \qquad
    \text{Wrong!}
  \end{equation*}

  This is \textbf{not} appending: here, the stack~\(s\) becomes an
  item of the resulting stack, but we want only the \textbf{contents}
  of~\(s\) to be pushed on top of~\(t\), not their container.

  \bigskip

  Analogy: You may think of an empty stack as a empty cardboard box,
  open topside, and every item pushed in it is like a book that you
  put on the topmost book, or at the bottom of the box if this is the
  first book.

\end{frame}

% ------------------------------------------------------------------------
%
\begin{frame}{Appending a stack}

  Let us set the abbreviations
  \begin{itemize}

    \item \(\el := \fun{nil}()\),

    \item \(\cons{x}{s} := \fun{cons}(x,s)\),

  \end{itemize}
  after the convention of the programming language \OCaml.

  \bigskip

  We can further abbreviate
  \begin{itemize}

    \item \(\cons{x_1}{\cons{x_2}{\cons{\dots}{\cons{x_n}{s}}}} :=
      \fun{cons}(x_1, \fun{cons}(x_2, \dots \fun{cons}(x_n,s)))\)

    \item \([x] := \cons{x}{\el}\)

    \item \([x_1; x_2, \dots; x_n] :=
          \cons{x_1}{\cons{x_2}{\cons{x_3}{\cons{\ldots}{\cons{x_n}{\el}}}}}\).

  \end{itemize}

  Our system now becomes a bit more legible:
  \begin{align*}
    \fun{cat}(        \el,t) &\xrightarrow{\alpha} t\\
    \fun{cat}(\cons{x}{s},t) &\xrightarrow{\smash{\beta}}
    \cons{x}{\fun{cat}(s,t)}
  \end{align*}

\end{frame}

% ------------------------------------------------------------------------
%
\begin{frame}{Pattern matching}

  \begin{itemize}

    \item Let us compute \(\fun{cat}([1;2],[3])\).

    \item We need to find the first rule that is matched by this call.

    \item Remember that the arguments must be values before you do
      that.

    \item Intuitively, this means that we compare the call with the
      patterns, that is, the left\hyp{}hand sides of the rules, in
      order.

    \item Here, rule~\(\alpha\) is not matched because \([1;2] \neq
      []\).

    \item Rule~\(\beta\) is matched because \([1;2] = 1::[2]\) and
      equals the sub\hyp{}pattern \(x::s\) if we assign \(x := 1\) and
      \(s := [2]\). Also, trivially, \(t := [3]\) works.

    \item After the rule~\(\beta\) is selected, the assignments are
      used to replace the variables in the right\hyp{}hand sides.

    \item The whole evaluation is
    \begin{equation*}
      \fun{cat}([1;2],[3])
      \xrightarrow{\beta}
      \cons{1}{\fun{cat}([2],[3])}
      \xrightarrow{\beta}
      \cons{1}{\cons{2}{\fun{cat}(\el,[3])}}
      \xrightarrow{\alpha}
                    [1;2;3].
    \end{equation*}
  \end{itemize}

\end{frame}

% ------------------------------------------------------------------------
%
\begin{frame}{Abstract Syntax Trees (ASTs)}

  \begin{itemize}

    \item Depending on the context, we may use the arborescent
      depiction of terms to bring to the fore certain aspects of a
      computation.

    \item For example, it may be interesting to show how parts of the
      output (the right\hyp{}hand side) are actually \textbf{shared}
      with the input (the left\hyp{}hand side).

    \item In other words, how much (new) data is created by a given
      rule.

    \item This concept supposes that terms reside in some sort of
      space and that they can be referred to from different other
      terms.

    \item This abstract space serves as a model of an interpreter's
      memory, called the \textbf{heap}.

  \end{itemize}

\end{frame}

% ------------------------------------------------------------------------
%
\begin{frame}{Data sharing}

  \begin{itemize}

    \item Consider for instance
      \begin{center}
        \includegraphics[bb=71 656 302 721]{cat_dag}
      \end{center}
      which is the same definition of `\fun{cat}' as given above:
      \begin{equation*}
        \fun{cat}(\el,t) \xrightarrow{\alpha} t
        \qquad\qquad\qquad
        \fun{cat}(\cons{x}{s},t) \xrightarrow{\beta}
        \cons{x}{\fun{cat}(s,t)}
      \end{equation*}

    \item The arrows on certain edges denote data sharing.

    \item When trees are used to visualise terms, they are called
      \textbf{abstract syntax trees} (AST).

    \item When some trees share subtrees, the whole \textbf{forest}
      (set of trees) is called a \textbf{directed acyclic graph}
      (DAG).

  \end{itemize}

\end{frame}

% ------------------------------------------------------------------------
%
\begin{frame}{Data sharing and evaluation}

  The evaluation of \(\fun{cat}([1;2],[3])\) is:
  \bigskip
  \begin{center}
    \includegraphics[bb=71 621 276 733]{cat123_push}
  \end{center}

  Note that each right-hand side shows only the rewritten call, not
  the whole term.

\end{frame}

% ------------------------------------------------------------------------
%
\begin{frame}{Data sharing and evaluation}

  Therefore, to reconstruct the value, we need to work leftwards and
  replace the rewritten call by the right-hand side:
  \begin{center}
    \includegraphics[bb=71 623 248 717]{cat123_pop0}
  \end{center}

  We do not show the calls that have been replaced by their values.

\end{frame}

% ------------------------------------------------------------------------
%
\begin{frame}{Data sharing and evaluation}

  \begin{center}
    \includegraphics[bb=71 621 212 715]{cat123_pop1}
    \qquad\quad
    \includegraphics[bb=71 621 212 715]{cat123_pop2}
  \end{center}

\end{frame}

% ------------------------------------------------------------------------
%
\begin{frame}{Call stack}

  We can see that the value of the second argument, [3], is shared
  with the value of the call, as well as the integers~1 and~2.

  \bigskip

  We saw how rewrites accumulated rightwards until the result had to
  be reconstructed by replacing leftwards the right-hand sides.

  \bigskip

  This dynamic is that of a stack, but a stack holding suspended
  function calls (that is, awaiting their values) and references to
  ASTs in the heap. This is the \textbf{call stack}.

  \bigskip

  The top of the stack (here, the rightmost abstract syntax tree) is
  the next term to be rewritten, unless the call on the left\hyp{}hand
  side is the original call, in which case the top is its value.

\end{frame}

% ------------------------------------------------------------------------
%
\begin{frame}{Garbage collection}

  \begin{itemize}

    \item In the abstract syntax trees (AST), the nodes with the calls
      that are replaced by their values have been immediately
      removed. In fact, they could have been removed later, in any
      order.

    \item The process which identifies and discards useless ASTs is
      called the \textbf{garbage collector} (GC).

    \item A term is useless if it cannot be accessed from anywhere in
      the call stack.

    \item Contrary to the heap, which can be thought of as an
      unordered set of trees (a forest of DAGs), the call stack is an
      ordered set that can be efficiently extended or reduced.

    \item The garbage collector can be conceived as a process
      interleaved with the process of evaluation.

    \item Like evaluators of functional programs, garbage collectors
      may also have different strategies (as to when dispose of
      useless data).
  \end{itemize}

\end{frame}

% ------------------------------------------------------------------------
%
\begin{frame}{Tail call optimisation}

  Let us define a function that sums integers in a non\hyp{}empty
  list:
  \begin{equation*}
    \fun{sum} \, [n] \xrightarrow{\alpha} n\qquad
    \fun{sum}(\cons{n}{s}) \xrightarrow{\beta} n + \fun{sum}(s)
  \end{equation*}
  where we introduced a notation `\(\fun{sum}[n]\)' instead of
  `\(\fun{sum}([n])\)'. We will omit the parentheses when there is
  exactly one argument and it is clear when it starts and ends.

  \bigskip

  The view of the system with maximum sharing is:
  \begin{center}
    \includegraphics{sum_dag}
  \end{center}

\end{frame}

% ------------------------------------------------------------------------
%
\begin{frame}{Tail call optimisation}

  The first phase of the evaluation of \(\fun{sum}([3;7;5])\) is:
  \begin{center}
    \includegraphics[bb=72 635 404 721,scale=0.99]{sum375_push}
  \end{center}

\end{frame}

% ------------------------------------------------------------------------
%
\begin{frame}{Tail call optimisation}

  The second phase (replacing the pending calls with their values) is
  \begin{center}
    \includegraphics[bb=72 644 380 721]{sum375_pop}
  \end{center}

\end{frame}

% ------------------------------------------------------------------------
%
\begin{frame}{Tail call optimisation}

  The same, with a garbage collector that collects after each
  replacement:
  \begin{center}
    \includegraphics[bb=71 644 378 721]{sum375_pop1}
  \end{center}

\end{frame}

% ------------------------------------------------------------------------
%
\begin{frame}{Tail call optimisation}

  Let us consider an alternate version of `\fun{sum}':
  \begin{equation*}
    \begin{array}{@{}r@{\;}l@{\;}l@{}}
      \fun{sum}_1(\cons{n}{t}) & \xrightarrow{\gamma}
      & \fun{sum}_0(n,t).\\
      \fun{sum}_0(n,\el) & \xrightarrow{\delta} & n;\\
      \fun{sum}_0(n,\cons{m}{s}) & \xrightarrow{\epsilon}
      & \fun{sum}_0(n+m,s).
    \end{array}
  \end{equation*}

  For \(\fun{sum}_1([3;7;5])\) we have the evaluation
  \begin{center}
    \includegraphics[bb=71 641 342 722]{sum0_375_0}
  \end{center}

\end{frame}

% ------------------------------------------------------------------------
%
\begin{frame}{Tail call optimisation}

  \begin{center}
    \includegraphics[bb=71 641 346 721]{sum0_375_1}
    \includegraphics[bb=71 639 384 721]{sum0_375_2}
  \end{center}

\end{frame}

% ------------------------------------------------------------------------
%
\begin{frame}{Tail call optimisation}

  Notice that, in the case of \fun{sum}$_0$, we found the result (15)
  at the end of the first phase (of growing the call stack).

  \bigskip

  Therefore, there is no need for a second phase, which replaces the
  pending calls by their values: we can discard all those calls at
  once after the first phase.

  \bigskip

  Or, to see it even more clearly, let us rewrite the same first phase
  and assume that the collector reclaims immediately the useless
  pending calls:
  \begin{center}
    \includegraphics[bb=71 640 326 721]{sum0_375tco0}
  \end{center}

\end{frame}

% ------------------------------------------------------------------------
%
\begin{frame}{Tail call optimisation}

  \begin{center}
    \includegraphics[bb=71 628 350 721]{sum0_375tco1}
    \includegraphics[bb=71 639 366 721]{sum0_375tco2}
  \end{center}

\end{frame}

% ------------------------------------------------------------------------
%
\begin{frame}{Tail call optimisation}

  As we see now, we can keep the size of the call stack constant, if
  we reclaim the top at each new rewrite.

  \bigskip

  This optimisation is called \textbf{tail call optimisation}.

  \bigskip

  Intuitively, a call in \textbf{tail position} is a call whose value
  is the value of the previously rewritten call.

  \bigskip

  For example, the call \(f(x)\) in \(1 + f(x)\) is not in tail
  position, because its value needs to be incremented.

  \bigskip

  Most compilers of functional languages detect and implement that
  optimisation that saves memory, including the \Clang compiler of
  \textsf{GCC}.

\end{frame}

% ------------------------------------------------------------------------
%
\begin{frame}{Inductive proofs}

Let us remark that \(\fun{cat}([1], [2,3,4]) \twoheadrightarrow
[1,2,3,4] \twoheadleftarrow \fun{cat}([1,2],
[3,4])\).

\bigskip

It is enlightening to create \emph{equivalence classes} of terms that
are joinable. These classes then define an equivalence
relationship~\((\equiv)\) as follows:
\begin{center}
  \(a \equiv b\) if there exists a unique value~\(v\) such that \(a
  \xrightarrow{\smash[t]{*}} v\) and \(b \xrightarrow{\smash{*}} v\).
\end{center}
For instance, \(\fun{cat}([1,2], [3,4]) \equiv \fun{cat}([1],
[2,3,4])\).

\bigskip

The relation (\(\equiv\)) is indeed an equivalence because it is
\begin{itemize}

  \item \emph{reflexive}: \(a \equiv a\);

  \item \emph{symmetric}: if \(a \equiv b\), then \(b \equiv a\);

  \item \emph{transitive}: if \(a \equiv b\) and \(b \equiv c\), then
    \(a \equiv c\).

\end{itemize}

\end{frame}

% ------------------------------------------------------------------------
%
\begin{frame}{Inductive proofs}

  Of some interest are the following facts.
  \begin{itemize}

    \item If \(x_1 \leftarrow x_2 \twoheadrightarrow x_3 \equiv x_4
      \rightarrow x_5 \twoheadleftarrow x_6\), then \(x_1 \equiv x_2
      \equiv x_3 \equiv x_4 \equiv x_5 \equiv x_6\).

    \item If \(f(x)\) and~\(f(y)\) have a value, then
      \[x \equiv y \Rightarrow f(x) \equiv f(y),\]
      that is, \((\equiv)\) is a \textbf{congruence relation}, as a
      consequence of confluence and termination.

  \end{itemize}

  If we want to prove equivalences with variables ranging over
  \textbf{infinite sets}, like \(\fun{cat}(s,\fun{cat}(t,u)) \equiv
  \fun{cat}(\fun{cat}(s,t),u)\), we need the \textbf{induction
    principle}.

\end{frame}

% ------------------------------------------------------------------------
%
\begin{frame}{Well-founded induction}

We define a \textbf{well\hyp{}founded order} on a set~\(A\) as being a
binary relation \((\succ)\) such that it is
\begin{itemize}

  \item irreflexive: \(\forall x \in A. \neg(x \succ x)\),

  \item transitive: \(\forall x, y, z \in A. x \succ y \,\wedge\, y
    \succ z \Rightarrow x \succ z\),

  \item and does not have any \textbf{infinite descending chains}, to
    wit, there is no \(a_0 \succ a_1 \succ a_2 \succ \dots\)

\end{itemize}
The \textbf{well\hyp{}founded induction principle} then states that, for
any predicate~\(\aleph\) on a set~\(A\),
\begin{center}
  \(\forall a \in A.\aleph(a)\) is implied by \(\forall a.(\forall b.a
  \succ b \Rightarrow \aleph(b)) \Rightarrow \aleph(a)\).
\end{center}
When \(A=\mathbb{N}\), this principle is called \textbf{mathematical
(complete) induction}.

\end{frame}

% ------------------------------------------------------------------------
%
\begin{frame}{Well-founded induction}

\textbf{Structural induction} is another particular case where \(t \succ
s\) holds if, and only if, \(s\)~is a proper \textbf{subterm} of~\(t\),
namely, the abstract syntax tree of~\(s\) is included in the tree
of~\(t\) and \(s \neq t\).

\bigskip

Sometimes, a restricted form of structural induction is enough. For
instance, we can define \(\cons{x}{s} \succ s\), for any term~\(x\)
and any stack~\(s \in S\). Both \(x\)~and~\(s\) are \textbf{immediate
subterms} of \(\cons{x}{s}\). (In other words, we dispensed with the
transitivity axiom.) Because there are no infinite descending chains,
any subset \(B \subseteq A\) contains minimal elements \(M \subseteq
B\), that is, there is no \(b \in B\) such that \(a \succ b\), if \(a
\in M\). So
\begin{enumerate}

  \item we have to prove \(\aleph(a)\) for all \(a \in M\), the
    so\hyp{}called \textbf{basis}: here, only \(\aleph(\el)\), as
    \(\el\)~is by construction the unique, minimal stacks, and,

  \item assuming the so\hyp{}called \textbf{induction hypothesis}
    \(\aleph(s)\), we must prove \(\aleph(\cons{x}{s})\) for
    all~\(x\).

\end{enumerate}

\end{frame}

% ------------------------------------------------------------------------
%
\begin{frame}{Termination}

  When defining our functional rewrite systems, we allowed programs to
  not terminate (see pages~\pageref{termination1}
  and~\pageref{termination2}).

  \bigskip

  We could actually have imposed some syntactic restrictions on
  recursive definitions in order to guarantee the termination of all
  functions. A well\hyp{}known class of such terminating functions
  makes exclusively use of a bridled form of recursion called
  \textbf{primitive recursion}.

  \bigskip

  Unfortunately, many useful functions do not fit, easily or at all,
  in this framework and, as a consequence, most functional languages
  leave to the programmers the responsibility to check the termination
  of their programs. For theoretical reasons, it is not possible to
  provide a general criterion for termination, but some rules exist
  that cover many usages.

\end{frame}

% ------------------------------------------------------------------------
%
\begin{frame}{Termination is undecidable}

  Here is an informal argument that shows that there is no predicate
  on functions that decides where their subject terminates on all
  inputs. By contradiction, let us assume that there exists~\(f\) such that
  \begin{equation*}
    f(g) =
    \begin{cases}
      \mathsf{true}  & \text{if $g$ terminates on all inputs},\\
      \mathsf{false} & \text{otherwise}.
    \end{cases}
  \end{equation*}
  Now let us define a function \(h\) as follows:
  \begin{equation*}
    h(g) =
    \begin{cases}
      \text{loops forever}            & \text{if $f(g)$},\\
      \text{terminates on all inputs} & \text{otherwise.}
    \end{cases}
  \end{equation*}
  In other words, \(h\) looks at \(f(g)\) and does the opposite of
  what \(f\) predicts about \(g\).

\end{frame}

% ------------------------------------------------------------------------
%
\begin{frame}{Termination is undecidable}

  Let us consider now \(h(h)\):
  \begin{itemize}

    \item If \(f(h)\) is true, then, by definition of~\(f\),
      \(h\)~terminates on all inputs, in particular
      \(h(h)\)~terminates. But, by definition of~\(h\), since \(f(h)\)
      holds, \(h(h)\) loops on all inputs, yielding a contradiction.

    \item If \(f(h)\) is false, then, by definition of~\(f\),
      \(h\)~does not terminate on all its inputs. But, by definition
      of~\(h\), since \(f(h)\) does not hold, \(h(h)\) terminates on
      all inputs, yielding another contradiction.

  \end{itemize}

\end{frame}

% ------------------------------------------------------------------------
%
\begin{frame}{Termination / Factorial}

  Here is the definition of the factorial function as a functional
  rewrite system:
  \begin{align*}
    \fun{fact} \, 0 & \rightarrow 1;\\
    \fun{fact} \, n & \rightarrow n \times \fun{fact}(n-1).
  \end{align*}

  \medskip

  Given a call \(\fun{fact}(n)\), with \(n \geqslant 0\), the rewrites
  yield a series of calls with decreasing arguments:
  \(\fun{fact}(n-1)\), \(\fun{fact}(n-2)\), etc. until
  \(\fun{fact}(0)\).

  \bigskip

  Termination here ensues from the \textbf{strictly decreasing
    arguments}:
  \begin{equation*}
    n > n-1 > n-2 > \dots > 0.
  \end{equation*}
  No infinite descent.

\end{frame}

% ------------------------------------------------------------------------
%
\begin{frame}{Termination / Ackermann's function}

Consider the following example where \(m,n \in \mathbb{N}\):
\begin{equation*}
\begin{array}{@{}r@{\;}l@{\;}l@{}}
\fun{ack}(0,n)     & \smashedrightarrow{\alpha} & n+1;\\
\fun{ack}(m+1,0)   & \smashedrightarrow{\beta} & \fun{ack}(m,1);\\
\fun{ack}(m+1,n+1) & \smashedrightarrow{\gamma}
                   & \fun{ack}(m,\fun{ack}(m+1,n)).
\end{array}
\end{equation*}
This is a simplified form of Ackermann's function, an early example of
a total computable function which is not primitive recursive. It makes
use of double recursion and two parameters to grow values as towers of
exponents, for example,
\begin{equation*}
  \fun{ack}(4,3) \twoheadrightarrow 2^{2^{65536}} - 3.
\end{equation*}
It is not obviously terminating, because if the first argument does
decrease, the second largely increases.

\end{frame}

% ------------------------------------------------------------------------
%
\begin{frame}{Termination / Ackermann's function}

Let us define a well\hyp{}founded ordering on pairs, called
\textbf{lexicographic order}. Let \((\succ_A)\) and \((\succ_B)\) be
well\hyp{}founded orders on the sets \(A\) and~\(B\). Then \({(\succ_{A
  \times B})}\) defined as follows on \(A \times B\) is
well\hyp{}founded:
\begin{equation*}
(a_0,b_0) \succ_{A \times B} (a_1,b_1) :\Leftrightarrow \text{\(a_0
    \succ_A a_1\) or (\(a_0 = a_1\) and \(b_0 \succ_B b_1\)).}
\end{equation*}
If \(A=B=\mathbb{N}\) then \((\succ_A) = (\succ_B) = (>)\). \\ For
example, \((\underline{1},7) \prec (\underline{2},3)\) and
\((1,\underline{7}) \succ (1,\underline{3})\).

\bigskip

To prove that \(\fun{ack}(m,n)\) terminates for all \(m,n \in
\mathbb{N}\), first, we must find a well\hyp{}founded order on the
calls \(\fun{ack}(m,n)\), that is, the calls must be totally ordered
without any infinite descending chain.

\bigskip

Here, a lexicographic order on \((m,n) \in \mathbb{N}^2\) extended to
\(\fun{ack}(m,n)\) works:
\begin{equation*}
\fun{ack}(a_0,b_0) \succ \fun{ack}(a_1,b_1) :\Leftrightarrow
\text{\(a_0 > a_1\) or (\(a_0 = a_1\) and \(b_0 > b_1\)).}
\end{equation*}

\end{frame}

% ------------------------------------------------------------------------
%
\begin{frame}{Termination / Ackermann's function}

  Clearly, \(\fun{ack}(0,0)\) is the minimum element.

  \bigskip

  Second, we must prove that \fun{ack} rewrites to smaller calls. We
  are only concerned with rules \(\beta\)~and~\(\gamma\). With the
  former, we have the inequality
  \begin{equation*}
    \fun{ack}(m+1,0) \succ \fun{ack}(m,1).
  \end{equation*}
  With rule~\(\gamma\), we have the following inequalities:
  \begin{equation*}
    \begin{array}{@{}r@{\;}l@{\;}l@{}}
      \fun{ack}(m+1,n+1) & \succ & \fun{ack}(m+1,n),\\
      \fun{ack}(m+1,n+1) & \succ & \fun{ack}(m,p),
    \end{array}
  \end{equation*}
  for all values~\(p\), in particular when \(\fun{ack}(m+1,n)
  \twoheadrightarrow p\).\hfill\(\Box\)

  \bigskip

  Implicitly, we used \[x \succ y \Rightarrow f(x) \succ f(y),\]
  that is, \textbf{evaluation is (strictly) order\hyp{}preserving}.

\end{frame}

% ------------------------------------------------------------------------
%
\begin{frame}{Termination / Dependency pairs and subterm order}

  In general, the approach often taken to prove termination consists
  in finding a well\hyp{}founded order \((\succ)\) on the rewritten
  terms, which is entailed by the rewrite relation \((\rightarrow)\),
  that is,
  \begin{equation*}
    x \rightarrow y  \Rightarrow  x \succ y.
  \end{equation*}

  \medskip

  As we have seen, we do not always can order the whole terms
  \(x\)~and~\(y\): only the calls in them to a given function are
  enough. These are called the \textbf{dependency pairs}. Calls to
  other functions can be assumed to be terminating.

  \bigskip

  With the factorial and Ackermann's function, we have seen how to
  deal with one or more integer arguments to these calls.

  \bigskip

  What about other data types, like stacks (sometimes called `lists')
  and trees?

\end{frame}

% ------------------------------------------------------------------------
%
\begin{frame}{Termination / Dependency pairs and subterm order}

  One well\hyp{}founded order for lists (stacks) is the
  \textbf{immediate subterm order}, satisfying
  \begin{equation*}
    \cons{x}{s} \succ s \quad \text{and} \quad \cons{x}{s} \succ x.
  \end{equation*}

  The immediate subterm order simply means that a list is `greater'
  than its first element and the rest of the list (called the
  \textbf{tail}). It is a special case of the \textbf{proper subterm
    order}, in which a structured value is `greater' than any of its
  components, wherever they are located.

  \bigskip

  These orders are convenient to prove the termination of calls to
  functions whose definitions are recursive in terms of substructures
  of their arguments: this is quite a common case.

  \bigskip

  For example, let us recall the function catenating a list to
  another:
  \begin{align*}
    \fun{cat}(        \el,t) &\xrightarrow{\alpha} t\\
    \fun{cat}(\cons{x}{s},t) &\xrightarrow{\smash{\beta}}
    \cons{x}{\fun{cat}(s,t)}
  \end{align*}
  We simply have: \(\cons{x}{s} \succ s \Rightarrow
  \fun{cat}(\cons{x}{s},t) \succ \fun{cat}(s,t).\)

\end{frame}

% ------------------------------------------------------------------------
%
\begin{frame}{Termination / Insertion sort}

  Let us consider \textbf{insertion sort}:
\begin{equation*}
\begin{array}{@{}r@{\;}l@{\;}lr@{\;}l@{\;}l@{}}
  \fun{isort}\,\el
& \smashedrightarrow{\alpha}
& \el;
& \fun{ins}(x, \cons{y}{s})
& \smashedrightarrow{\gamma}
& \cons{y}{\fun{ins}(x,s)}, \,\text{if \(x \succ y\)};\\
  \fun{isort}(\cons{x}{s})
& \smashedrightarrow{\beta}
& \fun{ins}(x,\fun{isort}\,s).
&  \fun{ins}(x,s)
& \smashedrightarrow{\delta}
& \cons{x}{s}.
\end{array}
\end{equation*}
Let us consider a short example:
\begin{equation*}
\begin{array}{r@{\;}l@{\;}l}
\fun{isort}\,[3,1,2]
& \smashedrightarrow{\beta} & \fun{ins}(3,\fun{isort}\,[1;2])\\
& \smashedrightarrow{\beta}
& \fun{ins}(3,\fun{ins}(1,\fun{isort}\,[2]))\\
& \smashedrightarrow{\beta}
& \fun{ins}(3,\fun{ins}(1,\fun{ins}(2,\fun{isort}\,\el)))\\
& \smashedrightarrow{\alpha}
& \fun{ins}(3,\fun{ins}(1,\fun{ins}(2,\el)))\\
& \smashedrightarrow{\delta}
& \fun{ins}(3,\fun{ins}(1,[2]))\\
& \smashedrightarrow{\delta}
& \fun{ins}(3,[1;2])\\
& \smashedrightarrow{\gamma}
& \cons{1}{\fun{ins}(3,[2])}\\
& \smashedrightarrow{\gamma}
& \cons{1}{\cons{2}{\fun{ins}(3,\el)}}\\
& \smashedrightarrow{\delta}
& \cons{1}{\cons{2}{\cons{3}{\el}}}\\
& = & [1;2;3].
\end{array}
\end{equation*}

\end{frame}

% ------------------------------------------------------------------------
%
\begin{frame}{Termination / Insertion sort}

  Informally, what soundness means is that, if some program
  terminates, then the result is what was expected. This property is
  called \textbf{partial correctness} when it is relevant to
  distinguish it from \textbf{total correctness}, which is partial
  correctness and termination.

  \bigskip

  Let us prove the termination of \fun{isort} by the method of
  dependency pairs. The pairs to order are the following:
  \begin{itemize}

    \item \((\fun{ins}(x, \cons{y}{s}), \fun{ins}(x,s))_\gamma\),
    \item \((\fun{isort}(\cons{x}{s}), \fun{isort}\,s)_\beta\) and
    \item \((\fun{isort}(\cons{x}{s}), \fun{ins}(x, \fun{isort}\,s))_\beta\).

  \end{itemize}

\end{frame}

% ------------------------------------------------------------------------
%
\begin{frame}{Termination / Insertion sort}

  By using the proper subterm relation on the second parameter of
  \fun{ins}, we order the first pair:
  \begin{equation*}
    \cons{y}{s} \succ s \Rightarrow \fun{ins}(x, \cons{y}{s}) \succ
    \fun{ins}(x,s).
  \end{equation*}
  This is enough to prove that \fun{ins} terminates. The second pair
  is similarly oriented:
  \begin{equation*}
    \cons{x}{s} \succ s \Rightarrow \fun{isort}(\cons{x}{s}) \succ
    \fun{isort}\,s.
  \end{equation*}

\end{frame}

% ------------------------------------------------------------------------
%
\begin{frame}{Termination / Insertion sort}

  The third pair is not worth considering after all, because we
  already know that \fun{ins} terminates, so the second pair is
  enough to entail the termination of \fun{isort}.

  \bigskip

  In other words, since \fun{ins} terminates, it can be considered,
  as far as termination analysis is concerned, as a data constructor,
  so the third pair becomes:
  \begin{equation*}
    \cons{x}{s} \succ s
    \Rightarrow
    \fun{isort}(\cons{x}{s}) \succ \fun{isort}\,s
    \Rightarrow
    \fun{isort}(\cons{x}{s}) \succ \underline{\fun{ins}}(x, \fun{isort}\,s).
  \end{equation*}
  where \(\fun{\ufun{ins}}\) stands for \fun{ins} considered as a
  \textbf{constructor}, that is, an undefined function.\hfill\(\Box\)

\end{frame}

% ------------------------------------------------------------------------
%
\begin{frame}{Termination / List flattening}

   Let us try a more complicated example:
   \begin{equation*}
     \begin{array}{r@{\;}l@{\;}l}
       \fun{flat}\,\el                   & \smashedrightarrow{\alpha}
       & \el\\
       \fun{flat}(\cons{\el}{t})         & \smashedrightarrow{\beta}
       & \fun{flat}\,t\\
       \fun{flat}(\cons{(\cons{x}{s})}{t}) & \smashedrightarrow{\gamma}
       & \fun{cat}(\fun{flat}(\cons{x}{s}),\fun{flat}\,t)\\
       \fun{flat}(\cons{y}{t})           & \smashedrightarrow{\delta}
       & \cons{y}{\fun{flat}\,t}
     \end{array}
   \end{equation*}

   Some evaluations:
   \begin{align*}
     \fun{flat}\,\el & \twoheadrightarrow \el\\
     \fun{flat}\,[\el;[\el]] & \twoheadrightarrow \el\\
     \fun{flat}\,[\el;[1;[2;\el];3];\el] & \twoheadrightarrow [1;2;3]
   \end{align*}

   What '\fun{flat}' does is flatten a list of lists arbitrarily
   deep.

\end{frame}

% ------------------------------------------------------------------------
%
\begin{frame}{Termination / List flattening}

  A detailed evaluation:
    \begin{equation*}
      \begin{array}{r@{\;}l@{\;}l}
        \fun{flat}\,[\el;[[1];2];3]
        & \xrightarrow{\beta}
        & \fun{flat}\,[[[1];2];3]\\
        & \smashedrightarrow{\gamma}
        & \fun{cat}(\underline{\fun{flat}\,[[1];2]},\fun{flat}\,[3])\\
        & \smashedrightarrow{\gamma}
        & \fun{cat}(\fun{cat}(\underline{\fun{flat}\,[1]},\fun{flat}\,[2]),
        \fun{flat}\,[3])\\
        & \smashedrightarrow{\delta}
        & \fun{cat}(\fun{cat}(\cons{1}{\underline{\fun{flat}\,\el}},
        \fun{flat}\,[2]), \fun{flat}\,[3])\\
        & \smashedrightarrow{\alpha}
        & \fun{cat}(\fun{cat}([1],\underline{\fun{flat}\,[2]}),
        \fun{flat}\,[3])\\
        & \smashedrightarrow{\delta}
        & \fun{cat}(\fun{cat}([1],\cons{2}{\underline{\fun{flat}\,\el}}),
        \fun{flat}\,[3])\\
        & \smashedrightarrow{\alpha}
        & \fun{cat}(\fun{cat}([1],[2]),\underline{\fun{flat}\,[3]})\\
        & \smashedrightarrow{\delta}
        & \fun{cat}(\fun{cat}([1],[2]),\cons{3}{\underline{\fun{flat}\,\el}})\\
        & \smashedrightarrow{\alpha}
        & \fun{cat}(\fun{cat}([1],[2]),[3])\\
        & \twoheadrightarrow & [1;2;3].
      \end{array}
    \end{equation*}

\end{frame}

% ------------------------------------------------------------------------
%
\begin{frame}{Termination / List flattening}

    The function \fun{cat} is independent of \fun{flat} and we
    proved its termination in isolation by using the proper subterm
    order on its first argument.

    \bigskip

    Assuming now that \fun{cat} terminates, let's prove the
    termination of \fun{flat}.

    \bigskip

    Because the recursive calls to \fun{flat} contain only (list)
    constructors, we can try to order their arguments.

    \bigskip

    The proper subterm order works:
    \begin{center}
    \begin{tabular}{r@{\;\,}c@{\;}ll}
      \(\cons{y}{t}\) & \(\succ\) & \(t\) & (rule~\(\delta\)
      and rule~\(\beta\) when \(y=\el\))\\
      \(\cons{(\cons{x}{s})}{t}\) & \(\succ\) & \(t\)
      & (rule~\(\gamma\))\\
      \(\cons{(\cons{x}{s})}{t}\) & \(\succ\) & \(\cons{x}{s}\)
      & (rule~\(\gamma\))
    \end{tabular}
    \end{center}

    \bigskip

    Termination ensues.\hfill\(\Box\)

\end{frame}

% ------------------------------------------------------------------------
%
\begin{frame}{Termination / List flattening}

  Let us consider now an alternative design for rule~\(\gamma\):
  \begin{equation*}
    \begin{array}{@{}r@{\;}c@{\;}l@{}}
      \fun{flat}\,\el & \xrightarrow{\alpha}
      & \el\\
      \fun{flat}(\cons{\el}{t}) & \smashedrightarrow{\beta}
      & \fun{flat}\,t\\
      \fun{flat}(\cons{(\cons{x}{s})}{t})
      & \smashedrightarrow{\gamma}
      & \fbox{\(\fun{flat}(\cons{x}{\cons{s}{t}})\)}\\
      \fun{flat}(\cons{y}{t})   & \smashedrightarrow{\delta}
      & \cons{y}{\fun{flat}\,t}
    \end{array}
  \end{equation*}

  \bigskip

  This amounts to perform a \textbf{right rotation} of the abstract
  syntax tree of the term:
  \begin{center}
    \includegraphics{flat_rotr}
  \end{center}

\end{frame}

% ------------------------------------------------------------------------
%
\begin{frame}{Termination / Measures}

  Here, the proper subterm order fails:
  \begin{equation*}
    \cons{(\cons{x}{s})}{t} \nsucc
    \cons{x}{(\cons{s}{t})} = \cons{x}{\cons{s}{t}}
  \end{equation*}
  because \(t \nsucc \cons{s}{t}\).

  \bigskip

  We need another method to prove termination: \textbf{measures}.

  \bigskip

  A measure \(\measure{\cdot}\) is a map from terms to a
  well\hyp{}ordered set (\(A, \succ\)), which is \textbf{monotonically
    decreasing} with respect to the rewrite relation:
  \begin{equation*}
    x \rightarrow y \Rightarrow \measure{x} \succ_A \measure{y}.
  \end{equation*}

  \medskip

  Actually, as earlier, we will only consider \textbf{dependency
    pairs}, that is, pairs of calls whose first component is the
  left\hyp{}hand side of a rule and the second components are the
  calls in the right\hyp{}hand side of the same rule.

\end{frame}

% ------------------------------------------------------------------------
%
\begin{frame}{Termination / Polynomial measures}

  The pairs are
  \begin{enumerate}

    \item \((\fun{flat}(\cons{\el}{t}), \fun{flat}\,t)_\beta\)

    \item \((\fun{flat}(\cons{(\cons{x}{s})}{t}),
      \fun{flat}(\cons{x}{\cons{s}{t}}))_\gamma\)

    \item \((\fun{flat}(\cons{y}{t}), \fun{flat}\,t)_\delta\)

  \end{enumerate}
  with \(y \not\in S\), that is, \(y\)~is not a list.

  \bigskip

  We can actually drop the function names, as all the pairs involve
  \fun{flat}.

  \bigskip

  A common class of measures are monotonically increasing embeddings
  into (\(\mathbb{N},>\)), so let us seek a measure satisfying:
  \begin{center}
    \begin{tabular}{r@{\;\,}c@{\;}ll}
      \(\measure{\cons{(\cons{x}{s})}{t}}\) & > &
      \(\measure{\cons{x}{\cons{s}{t}}}\)\\
      \(\measure{\cons{y}{t}}\) & > & \(\measure{t}\),
      & if \(y \not\in S\) or \(y=\el\).
    \end{tabular}
  \end{center}

\end{frame}

% ------------------------------------------------------------------------
%
\begin{frame}{Termination / Polynomial measures}

  For instance, let us set the following \textbf{polynomial measure}:
  \begin{center}
    \begin{tabular}{r@{\;\,}c@{\;\;}ll}
      \(\measure{\cons{x}{s}}\) & := & \(1 + 2\cdot\measure{x} +
      \measure{s}\)\\
      \(\measure{y}\) & := & \(0\), & if \(y \not\in S\) or
      \(y=\el\).
    \end{tabular}
  \end{center}

  We have
  \begin{align*}
    \measure{\cons{(\cons{x}{s})}{t}} &  =  3 +
    4\cdot\measure{x} + 2\cdot\measure{s} + \measure{t}\\
    \measure{\cons{x}{\cons{s}{t}}} &  =  2 +
    2\cdot\measure{x} + 2\cdot\measure{s} + \measure{t}
  \end{align*}

  Then
  \begin{center}
    \begin{tabular}{r@{\;\;}c@{\;\;}l}
      \(\measure{\cons{(\cons{x}{s})}{t}}\) & >
      & \(\measure{\cons{x}{\cons{s}{t}}}\)\\
      \(\measure{\cons{y}{t}} = 1 + \measure{t}\) & > &
      \(\measure{t}\)
    \end{tabular}
  \end{center}

  Since \(\measure{x} \in \mathbb{N}\), for all~\(x\), these
  inequalities entail the termination of \fun{flat} (no infinite
  descent).\hfill\(\Box\)

\end{frame}

% ------------------------------------------------------------------------
%
\begin{frame}{Termination / Dependency pairs in more details}

  Let's try to be a little bit more rigorous with dependency pair
  analysis.

  \bigskip

  Perhaps the first thing we have noticed is the use of conditions
  (guards) on rewrite rules. They are actually simple ways to make the
  systems more readable and more efficiently implemented in functional
  languages. In a way, they are akin to macros.

  \bigskip

  For example, consider the definitions of functions calculating the
  maximum and the minimum of two numbers:
  \begin{equation*}
    \begin{array}{@{}r@{\;}l@{\;}lr@{\;}l@{\;}l@{}}
      \fun{max}(0,y) & \rightarrow & y; & \fun{min}(0,y) & \rightarrow & 0;\\
      \fun{max}(x,0) & \rightarrow & x; & \fun{min}(x,0) & \rightarrow & 0;\\
      \fun{max}(x,y) & \rightarrow & 1 + \fun{max}(x\!-\!1,y\!-\!1).
      & \fun{min}(x,y) & \rightarrow & 1 + \fun{min}(x\!-\!1,y\!-\!1).
    \end{array}
  \end{equation*}
  Then insertion sort can be written (very inefficiently) as:
  \begin{equation*}
    \begin{array}{@{}r@{\;}l@{\;}l@{}}
      \fun{ins}(x,\el) & \rightarrow & [x];\\
      \fun{ins}(x,\cons{y}{s}) & \rightarrow &
      \cons{\fun{min}(x,y)}{\fun{ins}(\fun{max}(x,y),s)}.
    \end{array}
  \end{equation*}

\end{frame}

% ------------------------------------------------------------------------
%
\begin{frame}{Termination / Dependency pairs in more details}

  An equivalent but more compact definition would determine maximum
  and the minimum at the same time, like so:
\begin{equation*}
\begin{array}{r@{\;}c@{\;}l}
  \fun{isort}\,\el
& \smashedrightarrow{\alpha}
& \el;\\
  \fun{isort}(\cons{x}{s})
& \smashedrightarrow{\beta}
& \fun{ins}(x,\fun{isort}\,s).\\
  \fun{ins}(x, \el)
& \smashedrightarrow{\gamma}
& [x];\\
   \fun{ins}(x,\cons{y}{s})
& \smashedrightarrow{\delta}
& \fun{choose}(x, \cons{y}{s}, x, y).\\
  \fun{choose}(x, \cons{y}{s}, \fun{Succ}\,a, \fun{Succ}\,b)
& \smashedrightarrow{\epsilon}
& \fun{choose}(x, \cons{y}{s}, a, b);\\
  \fun{choose}(x, \cons{y}{s}, a, 0)
& \smashedrightarrow{\phi}
& \cons{y}{\fun{ins}(x,s)};\\
  \fun{choose}(x, \cons{y}{s}, 0, \fun{Succ}\,b)
& \smashedrightarrow{\psi}
& \cons{x}{\cons{y}{s}}.
\end{array}
\end{equation*}
Here, we use Peano's unary encoding of integers:
\begin{itemize}

  \item the successor function \fun{Succ}, so that \(\fun{Succ}(x)\)
    denotes canonically the value of \(x+1\); (We do not have to assume
    `-'.)

  \item Also, `0' here actually stands for \fun{Zero}().

\end{itemize}

\end{frame}

% ------------------------------------------------------------------------
%
\begin{frame}{Termination / Dependency pairs in more details}

  Instead of denoting the function symbols that are considered
  undefined in a dependency pair, as we did before, we are going to
  mark the defined ones (that is equivalent), with a sharp symbol
  (\#). This enables to write \textbf{the pairs as rewrite rules},
  which is more readable.

  \bigskip

  We have five of them:
  \begin{equation*}
    \begin{array}{r@{\;}c@{\;}ll}
        \fun{isort}^\#(\cons{x}{s})
      & \smashedrightarrow{a}
      & \fun{ins}^\#(x, \fun{isort}\, s), & \text{from \(\beta\)};\\
        \fun{isort}^\#(\cons{x}{s})
      & \smashedrightarrow{b}
      & \fun{isort}^\#\, s, & \text{from \(\beta\)};\\
        \fun{ins}^\#(x,\cons{y}{s})
      & \smashedrightarrow{c}
      & \fun{choose}^\#(x, \cons{y}{s}, x, y), & \text{from \(\delta\)};\\
        \fun{choose}^\#(x, \cons{y}{s}, \fun{Succ}\,a, \fun{Succ}\,b)
      & \smashedrightarrow{d}
      & \fun{choose}^\#(x, \cons{y}{s}, a, b), & \text{from \(\epsilon\)};\\
        \fun{choose}^\#(x, \cons{y}{s}, a, 0)
      & \smashedrightarrow{e}
      & \fun{ins}^\#(x,s), & \text{from \(\phi\)}.
    \end{array}
  \end{equation*}

\end{frame}

% ------------------------------------------------------------------------
%
\begin{frame}{Termination / Dependency pairs in more details}

  We can visualise better the pairs by drawing the dependency graph:
  each vertex is labelled by a pair (\(a\), \(b\), \(c\) etc.), and
  there is an arc \(i \rightarrow j\) if the right\hyp{}hand side of
  pair~\(i\) matches the left\hyp{}hand side of pair~\(j\):
  \begin{center}
    \begin{figure}
      \includegraphics[bb=80 675 268 705]{dp1}
    \end{figure}
  \end{center}
  That graph contains two \textbf{strongly\hyp{}connected components
    (SCCs)}: \(\{b\}\) and \(\{c, d, e\}\). We must show that
  \textbf{all cycles in each SCC map to a strictly decreasing
    measure.}

  \bigskip

  The first one is easy (it is a loop):
  \begin{equation*}
    \cons{x}{s} \succ s
    \Rightarrow
    \fun{isort}^\#(\cons{x}{s}) \smashedrightarrow{b} \fun{isort}^\#\,s
  \end{equation*}
  Another way to put it is that the projection of the pair~\(b\) on
  the parameter of \(\fun{isort}^\#\), that is \(\cons{x}{s}
  \rightarrow s\), if ordered, implies that the original pair is
  ordered too.

\end{frame}

% ------------------------------------------------------------------------
%
\begin{frame}{Termination / Dependency pairs in more details}

  We can annotate the graph with the ordering relation we found:
  \begin{center}
    \begin{figure}
      \includegraphics[bb=80 675 268 705]{dp2}
    \end{figure}
  \end{center}
  For the second and last SCC \(\{c, d, e\}\), if we project
  \(\fun{ins}^\#\) and \(\fun{choose}^\#\) on their second parameter,
  we get:
  \begin{equation*}
    \begin{array}{r@{\;}c@{\;}ll}
      \cons{y}{s} & \smashedrightarrow{c} & \cons{y}{s},\\
      \cons{y}{s} & \smashedrightarrow{d} & \cons{y}{s},\\
      \cons{y}{s} & \smashedrightarrow{e} & s.\\
    \end{array}
  \end{equation*}
  So all paths through~\(e\) strictly decreases the second parameter:
  \begin{center}
    \begin{figure}
      \includegraphics[bb=80 675 268 705]{dp3}
    \end{figure}
  \end{center}

\end{frame}

% ------------------------------------------------------------------------
%
\begin{frame}{Termination / Dependency pairs in more details}

  Now the SCC \(\{c, d, e\}\) reduces to the loop SCC \(\{d\}\).

  \bigskip

  Pair~\(d\) is easy to order: for example, we can project the
  components of the pair along the third parameter of
  \(\fun{choose}^\#\):
  \begin{equation*}
    \fun{Succ} \; a \xrightarrow{d} a.
  \end{equation*}
  So
  \begin{center}
    \begin{figure}
      \includegraphics[bb=80 675 268 705]{dp4}
    \end{figure}
  \end{center}
  (Remember that the immediate subterm ordering \((\succ)\) applies to
  the result of different projections.)

  \bigskip

  This entails that the system is terminating.\hfill$\Box$

\end{frame}

% ------------------------------------------------------------------------
%
\begin{frame}{Proving associativity of stack concatenation}

The definition of the concatenation of two stacks was:
\begin{equation*}
\fun{cat}(        \el,t) \xrightarrow{\smash{\alpha}} t;\qquad
\fun{cat}(\cons{x}{s},t) \xrightarrow{\smash{\beta}}
\cons{x}{\fun{cat}(s,t)}.
\end{equation*}
Let us prove the associativity of `\fun{cat}', which we express
formally as
\begin{equation*}
  \pred{CatAssoc}{s,t,u} \colon
\fun{cat}(s,\fun{cat}(t,u)) \equiv
\fun{cat}(\fun{cat}(s,t),u)
\end{equation*}
for all stack values \(s\), \(t\) and~\(u\).

\end{frame}

% ------------------------------------------------------------------------
%
\begin{frame}{Proving associativity of stack concatenation}

  We apply the well\hyp{}founded induction principle to the structure
  of~\(s\), so we must establish
  \begin{itemize}

    \item the basis \(\forall t,u \in S.\pred{CatAssoc}{\el,t,u}\), and the

    \item step \(\forall s,t,u \in S.\pred{CatAssoc}{s,t,u}
      \Rightarrow \forall x \in T.\pred{CatAssoc}{\cons{x}{s},t,u}\).

  \end{itemize}

  \bigskip

  Underlining the call being rewritten, the base case is:
  \begin{equation*}
    \ufun{cat}(\el,\fun{cat}(t,u))
    \xrightarrow{\smash{\alpha}} \fun{cat}(t,u)
    \xleftarrow{\smash{\alpha}} \fun{cat}(\ufun{cat}(\el,t),u).
  \end{equation*}

  Let us assume now \(\pred{CatAssoc}{s,t,u}\), called the
  \textbf{induction hypothesis}, and let us prove
  \(\pred{CatAssoc}{\cons{x}{s},t,u}\), for any term~\(x\).

\end{frame}

% ------------------------------------------------------------------------
%
\begin{frame}{Proving associativity of stack concatenation}

  The goal here is to use the rewrite system as an abstract machine to
  prove a property, with the timely help of the induction principle.

  \bigskip

  Precisely, we want to rewrite each side of the equivalence we wish
  to prove until we either find the same term (equality), or we use
  the induction hypothesis (equivalence).

  \bigskip

  Since we are going to work with equivalences, we should lift the
  call\hyp{}by\hyp{}value reduction strategy here, as it does not
  matter when arguments are evaluated, as long as their evaluation
  terminates and is confluent: we always assume these properties on
  the system hold when proving properties of the functions it defines.

\end{frame}

% ------------------------------------------------------------------------
%
\begin{frame}{Proving associativity of stack concatenation}

  We may start with the left\hyp{}hand side of
  \(\pred{CatAssoc}{\cons{x}{s},t,u}\), that is
  \(\fun{cat}(\cons{x}{s},\fun{cat}(t,u))\).

  \bigskip

  We notice that this is \textbf{not} a value, as the call contains a
  reducible call, \(\fun{cat}(t,u)\).

  \bigskip

  When using functions to compute, we match a call containing only
  values against the patterns of the rules.

  \bigskip

  Here, we would like to match a pattern against a pattern. This is
  more general than matching, because values are patterns, and is
  called \textbf{unification}. Given two patterns, we want to find
  whether we can find assignments to the variables of these so that,
  when they are substituted by their denotation, the two patterns
  become equal.

\end{frame}

% ------------------------------------------------------------------------
%
\begin{frame}{Proving associativity of stack concatenation}

 Here, we would like to unify
 \(\fun{cat}(\cons{x}{s},\fun{cat}(t,u))\) (called the \textbf{goal})
 and \(\fun{cat}(\cons{x}{s},t)\) (the pattern).

 \bigskip

The problem is that we have variables that are the same in the goal
and the pattern (\(x\), \(s\)~and~\(t\)).

\bigskip

The case of \(\cons{x}{s}\) and \(\cons{x}{s}\) is harmless, but the
unification of \(\fun{cat}(t,u)\) and~\(t\) is not.

\bigskip

That is why it is best to rename the variables in the pattern so they
are not found in the goal (this is the
\textbf{\(\alpha\)\hyp{}renaming} of \(\lambda\)\hyp{}calculus), and
then try to unify them.

\bigskip

To simplify things even more, we are going to rename \emph{all} the
variables in the definition of `\fun{cat}'.

\end{frame}

% ------------------------------------------------------------------------
%
\begin{frame}{Proving associativity of stack concatenation}

  Let us try for example the following renaming:
  \begin{align*}
    \fun{cat}(        \el,b) &\xrightarrow{\alpha} b\\
    \fun{cat}(\cons{i}{a},b) &\xrightarrow{\smash{\beta}}
    \cons{i}{\fun{cat}(a,b)}
  \end{align*}

  We want to unify the goal \(\fun{cat}(\cons{x}{s},\fun{cat}(t,u))\)
  with the pattern \(\fun{cat}(\cons{i}{a},b)\).

  \bigskip

  The unification yields the equations \(x = i\), \(s = a\) and
  \(\fun{cat}(t,u) = b\).

  \bigskip

  Then we apply the set of equations from the unification, so the
  right\hyp{}hand side of~\(\beta\), \(\cons{i}{\fun{cat}(a,b)}\),
  becomes \(\cons{x}{\fun{cat}(s,\fun{cat}(t,u))}\).

\end{frame}

% ------------------------------------------------------------------------
%
\begin{frame}{Proving associativity of stack concatenation}

  We have:
  \begin{equation*}
    \begin{array}{r@{\;}l@{\;}l@{\qquad}r@{}}
      \ufun{cat}(\cons{x}{s},\fun{cat}(t,u))
      & \xrightarrow{\smash\beta}
      & \cons{x}{\fun{cat}(s,\fun{cat}(t,u))}\\
      & \equiv
      & \cons{x}{\fun{cat}(\fun{cat}(s,t),u)}
      & \IH{\pred{CatAssoc}{s,t,u}}\\
      & \xleftarrow{\smash\beta}
      & \ufun{cat}(\cons{x}{\fun{cat}(s,t)},u)\\
      & \xleftarrow{\smash\beta}
      & \fun{cat}(\ufun{cat}(\cons{x}{s},t),u).
    \end{array}
  \end{equation*}

  Note the use of the inductive hypothesis \(\pred{CatAssoc}{s,t,u}\).

  \bigskip

  This derivation above establishes that
  \(\pred{CatAssoc}{\cons{x}{s},t,u}\) holds.

  \bigskip

  The induction principle entails that \(\forall s,t,u \in
  S.\pred{CatAssoc}{s,t,u}\).\hfill\(\Box\)

\end{frame}

% ------------------------------------------------------------------------
%
\begin{frame}{Unification}

  The best way to think about unification is to think of terms as
  trees, as we did page~\pageref{terms_as_trees}. A function is a
  subtree and a constant is a leaf.

  \bigskip

  The unification algorithm is then very simple:
  \begin{itemize}

    \item if the two terms are constants, they match if they are
      verbatim the same;

    \item if one of the terms is a variable, then they match;

    \item if the two terms are functions, they match if and only if
      \begin{itemize}

        \item they are the same (name),

        \item they have the same \textbf{arity} (number of arguments,
         that is, number of subterms, or number of subtrees),

        \item if they have arguments, the \(n\)-th argument of one
          term matches the \(n\)-th argument of the other.

  \end{itemize}

\end{itemize}

\end{frame}

% ------------------------------------------------------------------------
%
\begin{frame}{Unification / Example \#1}

  For example, let's unify the terms \(\fun{f}(\fun{g}(x,x),
  \fun{h}(y,y))\) and \(\fun{f}(\fun{g}(z), z)\).
  \begin{center}
    \includegraphics{tree-1-1}
  \end{center}
  They do not unify because the subterm \(\fun{g}(x,x)\) does not
  match \(g(z)\): the arity of \fun{g} is different.

  \bigskip

  Note how the terms are non-linear, even though we do not allow them
  in our systems. Indeed, unification is all about equalities, which
  non\hyp{}linearity is about. So let's allow it for exercises.

\end{frame}

% ------------------------------------------------------------------------
%
\begin{frame}{Unification / Example \#2}

  Consider now the terms \(\fun{f}(\fun{g}(x, x), \fun{h}(y, y))\) and
  \(\fun{f}(\fun{g}(\fun{h}(w, \fun{a}()), z), z))\):
  \begin{center}
    \includegraphics{tree-1-2}
  \end{center}

\end{frame}

% ------------------------------------------------------------------------
%
\begin{frame}{Unification / Example \#2}

  The steps are
\begin{align*}
  x &= \fun{h}(w, \fun{a}()) & x &= \fun{h}(w, \fun{a}())
  & \fun{h}(w, \fun{a}()) &= \fun{h}(y, y) & w &= y\\
x &= z      & x &= \fun{h}(y, y) & x      &= \fun{h}(y, y) & x &= \fun{h}(y, y)\\
z &= \fun{h}(y, y) & z &= x      & z      &= x      & y &= \fun{a}()\\
  &                &   &         &        &         & z &= x
\end{align*}
  \noindent Finally we obtain the \textbf{substitution}
\begin{align*}
w &= \fun{a}()\\
x &= \fun{h}(\fun{a}(), \fun{a}())\\
y &= \fun{a}()\\
z &= \fun{h}(\fun{a}(), \fun{a}())
\end{align*}

\end{frame}

% ------------------------------------------------------------------------
%
\begin{frame}{Unification / Example \#2}

  This answer means that if we take the two trees of the query and
  replace the variables according to the substitution, we get the
  exact same tree. Indeed, we get
  \begin{center}
    \includegraphics{tree-1-2u}
  \end{center}

\end{frame}

%% % ------------------------------------------------------------------------
%% %
%% \begin{frame}{Unification / Example \#3}

%%     Consider now the terms \(\fun{f}(\fun{g}(x,x), \fun{h}(\_,\_))\)
%%     and \(\fun{f}(\fun{g}(\fun{h}(w, a()),z), z)\), where the
%%     underscores denote \textbf{fresh variables} (that is, they are
%%     unknown but different from any other variable in the term):
%%     \begin{center}
%%       \includegraphics{tree-1-3}
%%     \end{center}

%% \end{frame}

%% % ------------------------------------------------------------------------
%% %
%% \begin{frame}{Unification / Example \#3}

%%   The steps are now
%%   \begin{align*}
%%     x &= \fun{h}(w,\fun{a}()) & x &= \fun{h}(w, \fun{a}())
%%       & x &= \fun{h}(w, \fun{a}()) & x &= \fun{h}(\alpha, \fun{a}())\\
%%     x &= z & x &= z & z &= x & z &= \fun{h}(\alpha, \fun{a}())\\
%%     \fun{h}(\alpha, \beta) &= z &
%%     \fun{h}(\alpha, \beta) &= \fun{h}(w, \fun{a}())
%%     & w &= \alpha & w &= \alpha\\
%%     & & & & & & \beta &= \fun{a}()
%%   \end{align*}

%%   Finally
%%   \begin{align*}
%%     \beta &= \alpha \\
%%     w &= \alpha\\
%%     x &= \fun{h}(\alpha, \fun{a}())\\
%%     z &= \fun{h}(\alpha, \fun{a}())
%%   \end{align*}

%% \end{frame}

%% % ------------------------------------------------------------------------
%% %
%% \begin{frame}{Unification / Example \#3}

%%   This means that if we replace the variables in the two initial
%%   trees, we get the same tree. Indeed, we get
%%   \begin{center}
%%     \includegraphics{tree-1-3u}
%%   \end{center}

%% \end{frame}

%% ------------------------------------------------------------------------
%
\begin{frame}{A mini functional language}

  Rewrite systems are a good mathematical and mental model for
  understanding functional languages.

  \bigskip

  It is time now to learn an actual functional language, so we can run
  programs. We are going to start with a mini language that is
  actually a subset of \OCaml, then grow it step by step.

  \bigskip

  Within rewrite systems, a `program' is simply a term that is reduced
  until reaching a stuck term. While this is a good theoretic
  framework to think about computation in general, it is not conducive
  to programming in the large, nor understanding the behaviour of
  programs.

  \bigskip

  We want programs to be made of different parts that can be reused in
  different contexts, for example, variables should be defined
  somewhere and reused at some another points.

\end{frame}

% ------------------------------------------------------------------------
%
\begin{frame}{Growing a language}

  This presentation is principled, that is, it relies on a systematic
  classification of the \textbf{syntactic constructs} and their
  \textbf{formal semantics}.

  \bigskip

  When we reach a certain complexity, we will drop the formal
  descriptions.

  \bigskip

  Both theoretical and pragmatic, our approach is complementary to
  those usually found in books for general audiences.

  \bigskip

  Let mini-ML be the name of the subset of \OCaml we are going to
  grow.

\end{frame}

% ------------------------------------------------------------------------
%
\begin{frame}{Appending a stack}

  The notation we introduced for stacks in the context of rewrite
  systems actually follows the \OCaml convention.

  \bigskip

  In the context of functional programming, stacks are often called
  \textbf{lists}.

  \bigskip

  Here are different ways to express the same list in \OCaml:
  \begin{equation*}
    1::(2::3::(4::\el)) = 1::2::3::4::\el = [1;2;3;4].
  \end{equation*}

  Before starting, let us see how a familiar function defined with a
  rewrite system looks like in \OCaml.

  \bigskip

  Let us recall the appending of a stack:
  \begin{align*}
    \fun{cat}(        \el,t) &\smashedrightarrow{\alpha} t;\\
    \fun{cat}(\cons{x}{s},t) &\smashedrightarrow{\beta} \cons{x}{\fun{cat}(s,t)}.
  \end{align*}

\end{frame}

% ------------------------------------------------------------------------
%
\begin{frame}[fragile]{Appending a stack}

  In \OCaml, the function name is not repeated (once for each rule).

  \bigskip

  Instead, it is introduced once by a \textbf{binder} called `\Xlet',
  like so:
  \vspace*{-3mm}
\begin{alltt}
\small
\textbf{let} cat = \textbf{function} ...
\end{alltt}
  The complete definition in \OCaml is:
  \vspace*{-3mm}
\begin{alltt}
\small
\textbf{let} \textbf{rec} cat = \textbf{function}
\ \ \ [], t \(\rightarrow\) t
\(\mid\) \ttcons{x}{s}, t \(\rightarrow\) \ttcons{x}{\fun{cat}(s,t)}
\end{alltt}

   Note the addition of the keyword \Xrec: this is meant so the call
   \(\fun{cat}(s,t)\) indeed refers to the function being defined.

   \bigskip

   With rewrite systems, recursion is not an issue: any call matching
   a pattern can be rewritten.

   \bigskip

   Note also how the two rules are introduced by the keyword
   \Xfunction and are separated by a vertical line.

\end{frame}

% ------------------------------------------------------------------------
%
\begin{frame}{Sentences}

  An \OCaml program is a sequence of sentences. A \textbf{sentence},
  or \textbf{global definition}, is defined by the following cases,
  where~\(e\) denotes an expression, \(x\)~is a variable, and `\Xlet'
  and~`\Xrec' are keywords:

  \medskip

  \begin{tabular}{rll}
    $\bullet$ & \emph{global definition}
    & \phrase{$\Xlet \; x \; = \; e$}\\
    $\bullet$ & \emph{global recursive definition}
    & \phrase{$\Xlet \; \Xrec \; x \; = \; e$}
  \end{tabular}

  \medskip

  Global definitions can be optionally followed by two semicolons:
  these are necessary when inputting the sentences in the toplevel
  loop (prompted after running the command \texttt{\small ocaml} from
  a Unix shell).

  \bigskip

  Note that the keyword `\Xrec' is necessary to enable recursion. (We
  will see later how that works precisely.)

  \bigskip

  For~\(e\) to be an expression means that \(e\)~denotes a part of a
  sentence which is classified as an expression according to its
  \textbf{syntax}.

\end{frame}

% ------------------------------------------------------------------------
%
\begin{frame}[fragile]{Metavariables}

  \begin{itemize}

    \item We say that \(e\)~is a \textbf{metavariable} because, being
      a name, it is a variable, but that variable does not belong to
      the language being described (the \textbf{object language},
      here, \OCaml) and, instead, exists only in the descriptive
      language (the \textbf{metalanguage}).

    \item In the following
      \begin{center}
        \phrase{$\Xlet \; x \; = \; e$}
      \end{center}
      the metavariable~\(x\) (mind the italics) denotes an infinite
      set of \OCaml variables and we must not confuse it with the
      \OCaml variable~\ident{x} (no italics). Similarly, the
      metavariable \(e\)~denotes an infinity of expressions, as in the
      following:
      \vspace*{-3mm}
\begin{alltt}
\small
\textbf{let} e = 3
\textbf{let} x = e
\textbf{let} a = x + e
\end{alltt}

  \end{itemize}

\end{frame}

% ------------------------------------------------------------------------
%
\begin{frame}{Expressions}

  \begin{itemize}

    \item An expression \(e\)~is \textbf{recursively defined by} the
      following \textbf{cases}:

      \smallskip

      \begin{tabular}{@{\,}r@{\,\,}ll}
        $\bullet$
        & \emph{variable}
        & my\_var, x, y, z, etc. \\
        $\bullet$
        & \emph{unit or integer constant}
        & \unit \ or \textsf{0} or \textsf{1} or \textsf{2}, etc.\\
        $\bullet$
        & \emph{function} (or \emph{abstraction})
        & $\Xfun \; x \rightarrow e$\\
        $\bullet$
        & \emph{call} (or \emph{application})
        & $e_1 \; e_2$ \\
        $\bullet$
        & \emph{arithmetic operator}
        & \lpar\texttt{+}\rpar \ \lpar\texttt{-}\rpar \ \lpar\texttt{/}\rpar
        \ \lpar\texttt{*}\rpar\\
        $\bullet$
        & \emph{arithmetic operation}
        & $e_1$ \texttt{+} $e_2$ or $e_1$ \texttt{-} $e_2$,
        etc.\\
        $\bullet$
        & \emph{parenthesis}
        & $\lpar{e}\rpar$\\
        $\bullet$
        & \emph{local definition}
        & $\Xlet \; x = e_1 \; \Xin \; e_2$
      \end{tabular}

    \item Note how the keywords `\Xrec' and `\Xfunction', as well as
      the definition of multiple cases (as seen in the definition of
      `\fun{cat}' earlier) is not found above: this is because we want
      to focus on other aspects of \OCaml first, like
      higher\hyp{}order functions.

    \item Instead, we have `\Xfun', which enables only one case made
      of a variable (more generally, irrefutable patterns).

  \end{itemize}

\end{frame}

% ------------------------------------------------------------------------
%
\begin{frame}[fragile]{Examples of expressions}

  \begin{itemize}

    \item Note that what we print nicely as `\(\rightarrow\)' is
      actually written `\texttt{->}' in the source code.

    \item Here are examples of \OCaml sentences making up a program,
      step by step.

    \item First, a numerical constant:
      \vspace*{-3mm}
\begin{alltt}
\small
\textbf{let} x = 0
\end{alltt}

    \item Note that constants can be defined by means of a rewrite
      system as functions taking the empty tuple:
      \begin{equation*}
        \fun{x}() \rightarrow 0
      \end{equation*}

    \item The next sentence is
      \vspace*{-3mm}
\begin{alltt}
\small
\textbf{let} id = \textbf{fun} x \(\rightarrow\) x
\end{alltt}

    \item This is equivalent to the rewrite system
      \begin{equation*}
        \fun{id}(x) \rightarrow x
      \end{equation*}

  \end{itemize}

\end{frame}

% ------------------------------------------------------------------------
%
\begin{frame}[fragile]{Examples of expressions}

  \begin{itemize}

    \item Next we have another constant being defined by binding a
      variable to the value of a call:
      \vspace*{-3mm}
\begin{alltt}
\small
\textbf{let} y = 2 \textbf{in} id y
\end{alltt}
      This is equivalent to the call \texttt{\small id(2)}.

    \item Now
      \vspace*{-3mm}
\begin{alltt}
\small
\textbf{let} x = (\textbf{fun} x \(\rightarrow\) \textbf{fun} y \(\rightarrow\) x + y) 1 2
\end{alltt}
      is equivalent to
      \vspace*{-3mm}
\begin{alltt}
\small
\textbf{let} x = 1 + 2
\end{alltt}
      (We will see why later.)

    \item Finally
      \vspace*{-3mm}
\begin{alltt}
\small
\textbf{let} z = x + 1
\end{alltt}
       is obvious.

  \end{itemize}

\end{frame}

% ------------------------------------------------------------------------
%
\begin{frame}{Functions}

  \begin{itemize}

    \item Variables must start with a lowercase character.

    \item The arrow is right-associative, so the expression
      \begin{equation*}
        \Xfun \; x_1 \rightarrow \Xfun \; x_2
        \rightarrow \ldots
        \rightarrow \Xfun \; x_n \rightarrow e
      \end{equation*}
      is equivalent to \(\Xfun \; x_1 \rightarrow (\Xfun
      \; x_2 \rightarrow (\ldots \rightarrow
      (\Xfun \; x_n \rightarrow e)) \; \ldots)\).

    \item Function calls are left\hyp{}associative, so
      \begin{equation*}
        e_1 \; e_2 \; e_3 \; \ldots \; e_n \;=\; (((\ldots
        (e_1 \; e_2) \; e_3) \ldots) \; e_n).
      \end{equation*}

    \item Function calls have higher priority than operator calls, for
      example,
      \begin{equation*}
        f \; 3 + 4 \; = \; (f \; 3) + 4 \; = \; f(3) + 4.
      \end{equation*}

    \item Operator calls have higher syntactic priority than
      abstractions, for instance,
      \begin{equation*}
        \Xfun \; x \rightarrow x + y \quad\text{is the same as}\quad
        \Xfun \; x \rightarrow (x + y).
      \end{equation*}

  \end{itemize}

\end{frame}

% ------------------------------------------------------------------------
%
\begin{frame}[fragile]{Function calls}

  Note the difference between \texttt{\small f x y} and \texttt{\small
    f(x,y)}.

  \bigskip

  The former is to be understood as \texttt{\small (f(x))(y)} which
  means that there are \emph{two} calls:
  \begin{enumerate}

  \item First, the call \texttt{\small f(x)}
    is evaluated. Let us assume that its value is \texttt{\small g}.

  \item Second, the call \texttt{\small g(y)} is evaluated.

  \end{enumerate}

  In other words, the first case is equivalent to
  \vspace*{-3mm}
\begin{alltt}
\small
\textbf{let} g = f x \textbf{in} g y
\end{alltt}

\end{frame}

% ------------------------------------------------------------------------
%
\begin{frame}{Abstract syntax}

  A graphical representation of programs (here, expressions) as trees
  proves very handy when studying their semantics (meaning):

  \begin{center}
    \begin{tabular}{lc}
      \toprule
      \multicolumn{1}{c}{Expression}
        & \multicolumn{1}{c}{Tree}\\
        \toprule
        $x$
        & $x$\\
        $\Xfun \; x \rightarrow e$
        & \includegraphics[bb=148 638 171 662]{fun_tree}\\
        $e_1 \; e_2$
        & \includegraphics[bb=148 637 177 663]{app_tree}\\
        \bottomrule
    \end{tabular}
  \end{center}

  Note that we needed some symbol for the root of the application tree
  (\texttt{\small \$}). Traditionally, that symbol is (\texttt{\small
    @}), but \OCaml uses it already for something else.

  \bigskip

  When trees represent programs, they are called \textbf{abstract
    syntax trees} (AST).

\end{frame}

% ------------------------------------------------------------------------
%
\begin{frame}{Abstract syntax}

  Intuitively, the method for constructing abstract syntax trees
  consists in fully parenthesising the expression. Each parenthesis
  corresponds to a subexpression and each subexpression corresponds to
  a subtree.

  \bigskip

  The tree is built from the root, down to the leaves, by traversing
  subexpressions from the outermost to the innermost, that is, the
  most embedded ones. For instance, consider
  \begin{center}
    \includegraphics[bb=148 617 233 660]{arith_tree1}
    \qquad
    \includegraphics[bb=148 599 246 662]{arith_tree2}
  \end{center}
  \centerline{\texttt{\small ((1+2)*(5/1)) \quad\qquad \texttt{\small \textbf{let} x = 1 \textbf{in} (1+2)*5/1}}}

\end{frame}

% ------------------------------------------------------------------------
%
\begin{frame}[fragile]{Abstract syntax}

  Another example:
    \vspace*{-3mm}
\begin{alltt}
\small
\textbf{let} x = 1 \textbf{in} ((\textbf{let} x = 2 \textbf{in} x) + x)
\end{alltt}
  corresponds to the abstract syntax tree
  \begin{center}
    \includegraphics[bb=148 599 221 662]{arith_tree3}
  \end{center}

  This example features two local definitions of `x': is one the
  redefinition of the other? Is one hiding the other?

  \bigskip

  There are also two uses of `x': which use corresponds to which
  definition?

\end{frame}

% ------------------------------------------------------------------------
%
\begin{frame}{Bindings, environments, scope}

  To answer theses questions, we must first understand
  \begin{itemize}

    \item what a \textbf{binding} is,

    \item what the \textbf{scope} of a binding is, and

    \item what an \textbf{environment} is.
  \end{itemize}

  The expression we called \textbf{local definition} earlier:
  \begin{equation*}
    \Xlet \; x = e_1 \; \Xin \; e_2
  \end{equation*}
  is evaluated as follows:
  \begin{enumerate}

    \item The expression \(e_1\)~is evaluated,

    \item assuming it has a value~\(v_1\), then \(v_1\) is
      \textbf{bound} to the variable~\(x\),

    \item the expression \(e_2\)~is evaluated under the assumption
      that \(x\)~is bound to~\(v_1\).

  \end{enumerate}

\end{frame}

% ------------------------------------------------------------------------
%
\begin{frame}{Bindings, environments, scope}

  A local definition
  \begin{equation*}
    \Xlet \; x = e_1 \; \Xin \; e_2
  \end{equation*}
  associates the value~\(v_1\) of expression~\(e_1\) to the variable
  denoted by~\(x\), making up a \textbf{binding} noted \(x \; \mapsto
  \; v_1\), and then proceeds to evaluate~\(e_2\) with this binding
  available.

  \bigskip

  The binding \(x \mapsto v_1\) is not usable outside the evaluation
  of~\(e_2\). We say that the \textbf{scope} of the definition is made
  of the subtrees of the abstract syntax tree corresponding to~\(e_2\)
  where the binding of that definition is usable (`visible').

  \bigskip

  A mapping from variables to values is called an
  \textbf{environment}. Let us note \(\rho_\varnothing\) the empty
  environment. We write
  \begin{equation*}
    (x \mapsto v) \oplus \rho \quad\text{or, simply,}\quad x \mapsto v
    \oplus \rho,
  \end{equation*}
  to mean that the binding \(x \mapsto v\) was added to the
  environment~\(\rho\).

\end{frame}

% ------------------------------------------------------------------------
%
\begin{frame}[fragile]{Bindings, environments, scope}

  Environments are functions and are updated as follows:
  \begin{equation}
    (x \mapsto v \oplus \rho )(y) \triangleq
    \begin{cases}
      v         & \text{if} \; x = y,\\
      \sigma(y) & \text{otherwise.}
    \end{cases}
  \end{equation}
  with \(\forall x.\rho_\varnothing(x) = \bot\) and \(\bot\) is
  interpreted as an \textbf{error}. Errors are special
  values. Consider again the program
  \vspace*{-3mm}
\begin{alltt}
\small
\textbf{let} x = 0 \textbf{in}
\textbf{let} id = \textbf{fun} x \(\rightarrow\) x \textbf{in}
\textbf{let} y = id x \textbf{in}
\textbf{let} x = (\textbf{fun} x \(\rightarrow\) \textbf{fun} y \(\rightarrow\) x + y) 1 2
\textbf{in} x + 1
\end{alltt}

  The variable~\texttt{\small x} in the expression \texttt{\small x+1}
  denotes the value of the expression bound to the
  variable~\texttt{\small x} in the previous line, not the first.

  \bigskip

  Bindings are ordered in the environment \textbf{by the order of
    their definitions}. Let us follow the evolution of the environment.

\end{frame}

% ------------------------------------------------------------------------
%
\begin{frame}{Bindings, environments, scope}

  \begin{enumerate}

    \item The environment is initially empty: \(\rho_\varnothing\);

    \item \texttt{\small \textbf{let} x = 0 \textbf{in}
      \ldots}\\ yields $(\ident{x} \mapsto 0) \oplus
      \rho_\varnothing$;

    \item \texttt{\small \textbf{let} id = \textbf{fun} x
      $\rightarrow$ x \textbf{in} \ldots}\\ yields $(\ident{id}
      \mapsto \Xfun \; \ident{x} \rightarrow \ident{x}) \oplus
      (\ident{x} \mapsto 0) \oplus \rho_\varnothing$;

    \item \texttt{\small \textbf{let} y = id x \textbf{in}
      \ldots}\\ yields $(\ident{y} \mapsto 0) \oplus (\ident{id}
      \mapsto \Xfun \; \ident{x} \rightarrow \ident{x}) \oplus
      (\ident{x} \mapsto 0) \oplus \rho_\varnothing$;

    \item \texttt{\small \textbf{let} x = \ldots}\\ yields
      $(\underline{\ident{x} \mapsto 3}) \oplus (\ident{y}
      \mapsto 0) \oplus (\ident{id} \mapsto \Xfun \;
      \ident{x} \rightarrow \ident{x}) \oplus (\ident{x} \mapsto
      0) \oplus \rho_\varnothing$.

    \item The expression \textsf{\small x+1} is evaluated in \(3+1=4\)
      because \(\textsf{x} \mapsto 3\) (the last binding)
      \textbf{hides} \(\textsf{x} \mapsto 0\) (the first
      binding).

    \item In other words, the first definition of~\texttt{\small x} was
      \textbf{out of scope} in \texttt{\small x+1}.

  \end{enumerate}

\end{frame}

% ------------------------------------------------------------------------
%
\begin{frame}{Free variables}

  \begin{itemize}

    \item If a variable~\(x\) is not bound to a value in an
      environment~\(\rho\), we say that \(x\)~is \textbf{free}
      in~\(\rho\).

    \item By definition, a variable is \textbf{free} in an expression
      if it is not bound by a `\Xlet' or a `\Xfun'.

    \item We can understand this concept on a graphic representation
      of the abstract syntax tree.

    \item From each variable occurrence, we move up, towards the
      root.

    \item If we encounter a `\Xlet' that binds that variable (in the
      left child), we create an oriented edge from the variable to
      that `\Xlet', and we stop: the variable is bound.

    \item Otherwise, if we reach the root (no such `\Xlet' has been
      found), then the variable is free.

  \end{itemize}

\end{frame}

% ------------------------------------------------------------------------
%
\begin{frame}{Free variables}

  \begin{figure}
    \begin{center}
      \includegraphics[bb=148 599 221 662]{free_var1}
      \caption{\textbf{let} x = 1 \textbf{in}
      ((\textbf{let} x = 2 \textbf{in} x) + x)}
    \end{center}
  \end{figure}

  \begin{figure}
    \begin{center}
      \includegraphics[bb=148 582 211 662]{free_var2}
    \end{center}
    \caption{\textbf{fun} y \(\rightarrow\) x + (\textbf{fun} x
      \(\rightarrow\) x) y}
  \end{figure}

\end{frame}

% ------------------------------------------------------------------------
%
\begin{frame}{Free variables}

  A similar situation arises with functions: `\Xfun' is a binder, just
  like `\Xlet'.

  \bigskip

  In \Xfun $x \rightarrow e$, the variable~\(x\) (called a
  \textbf{parameter}) may shadow (hide) another variable~\(x\) defined
  on the path to the root.

  \bigskip

  See previous page for two examples.

  \bigskip

  We note \(\FV{e}\) the set of the variables free in~\(e\).

\end{frame}

% ------------------------------------------------------------------------
%
\begin{frame}{Free variables}

  It is easy to formally define \(\mathcal{F}\) by case and
  recursively:
  \begin{align*}
    \FV{x} &= \{x\}, \quad\text{where $x$ is a  variable}\\
    \FV{c} &= \varnothing, \quad\text{where $c$ is a constant}\\
    \FV{\Xfun \; x \rightarrow e} &=
    \FV{e} \backslash \{x\}\\
    \FV{e_1 \; e_2} &= \FV{e_1} \;\cup\; \FV{e_2}\\
    \FV{\Xlet \; x  = e_1 \; \Xin \; e_2} &=
    \FV{e_1} \;\cup\; \FV{e_2} \backslash \{x\}\\
    \FV{e_1 \; \texttt{+} \; e_2} &= \FV{e_1} \;\cup\; \FV{e_2}
  \end{align*}

  Now, instead of drawing abstract syntax trees, let us use these
  equations to compute the free variables of our previous examples.

\end{frame}

% ------------------------------------------------------------------------
%
\begin{frame}{Free variables}

  \begin{equation*}
    \begin{array}{lcl}
      \multicolumn{3}{l}{
        \FV{\Xlet \; \textsf{x} = 1 \;
          \Xin \; (\Xlet \; \textsf{x} = 2 \;
          \Xin \; \textsf{x}) \; + \; \textsf{x}}} \\
      \phantom{XXX} &=& \FV{1} \cup \FV{(\Xlet \; \textsf{x}
        = 2 \; \Xin \; \textsf{x}) +
        \textsf{x}} \backslash \{\textsf{x}\}\\
      \phantom{XXX} &=& \varnothing \cup (\FV{\Xlet \;
        \textsf{x} = 2 \; \Xin \;
        \textsf{x}} \cup \FV{\textsf{x}}) \backslash \{\textsf{x}\}\\
      \phantom{XXX} &=& (\FV{2} \cup \FV{\textsf{x}} \backslash
      \{\textsf{x}\}  \cup \FV{\textsf{x}}) \backslash
      \{\textsf{x}\}\\
      \phantom{XXX} &=& \FV{\textsf{x}} \backslash \{\textsf{x}\}\\
      \phantom{XXX} &=& \varnothing.
    \end{array}
  \end{equation*}

  \bigskip

  \begin{equation*}
    \begin{array}{lcl}
      \multicolumn{3}{l}{ \FV{\Xfun \; \textsf{y} \rightarrow
          \textsf{x} + (\Xfun \; \textsf{x}
          \rightarrow \textsf{x}) \; \textsf{y}}}\\ \phantom{X}
      &=& \FV{\textsf{x} + (\Xfun \; \textsf{x} \rightarrow
        \textsf{x}) \; \textsf{y}} \backslash
      \{\textsf{y}\}\\ \phantom{X} &=& (\FV{\textsf{x}} \cup
      \FV{(\Xfun \; \textsf{x} \rightarrow \textsf{x})
        \; \textsf{y}}) \backslash \{\textsf{y}\}\\ \phantom{X} &=&
      (\{\textsf{x}\} \cup \FV{\Xfun \; \textsf{x}
        \rightarrow \textsf{x}} \cup \FV{\textsf{y}})
      \backslash \{\textsf{y}\}\\ \phantom{X} &=& (\{\textsf{x}\}
      \cup \FV{\textsf{x}} \backslash \{\textsf{x}\} \cup
      \{\textsf{y}\}) \backslash \{\textsf{y}\}\\ \phantom{XXX} &=&
      \{\textsf{x}, \textsf{y}\} \backslash
      \{\textsf{y}\}\\ \phantom{X} &=& \{\textsf{x}\}.
    \end{array}
  \end{equation*}

\end{frame}

% ------------------------------------------------------------------------
%
\begin{frame}{Closed expressions}

  \begin{itemize}

    \item A \textbf{closed} expression is an expression without free
      variables.

    \item That is, an expression~\(e\) is closed if, and only if,
      \(\FV{e} = \varnothing\).

    \item By extension, we say that~\(e\) is \textbf{closed in an
      environment}~\(\rho\) if \(\FV{e} \subseteq \dom{\rho}\) (the
      free variables are bound in~\(\rho\)).

    \item Only a closed expression can be evaluated (executed). That
      is why the first static analysis performed by compilers consists
      in determining the variables which are free in expressions: if
      the program is not closed, it is rejected.

    \item In the case of
      \begin{center}
        \Xfun y $\rightarrow$ x + (\Xfun x $\rightarrow$ x) y,
      \end{center}
      the \OCaml compiler prints out
      \begin{center}
        \texttt{\small \textbf{Unbound value x}}
      \end{center}
      and stops. It is useful that this open expression be rejected at
      compile\hyp{}time and does not cause a run\hyp{}time error.

  \end{itemize}

\end{frame}

% ------------------------------------------------------------------------
%
\begin{frame}{Evaluation}

  \textbf{Evaluation} is a partial function from expressions to
  values. An \textbf{interpreter} is a program implementing an
  evaluation. The partiality models the fact that an evaluation may
  not terminate or end in an error.

  \bigskip

  With rewrite systems, values are defined as stuck expressions, that
  is, expressions that cannot be rewritten further.

  \bigskip

  Here, since we want to define expressions (terms) whose value we can
  reuse somewhere else (in another term), we need the notion of
  \textbf{environment}.

  \bigskip

  That implies that a formal account of evaluation (the equivalent for
  programs of rewriting terms) needs to take into account
  environments.

\end{frame}

% ------------------------------------------------------------------------
%
\begin{frame}{Semantic values}

  This goes even further because we want Mini-ML to be
  \textbf{higher\hyp{}order}, that is, we want functions to be values
  and this begets the question of the meaning of the free variables in
  the function bodies: after passing around a function as argument to
  other functions, what are the values of those variables when the
  threaded function is finally applied?

  \bigskip

  The most common answer is \textbf{lexical (or static) scoping},
  whereby the free variables in a body remain bound to their values at
  the point the function was defined.

  \bigskip

  This means that values corresponding to functions have to carry
  around the environment at the point of their definition.

  \bigskip

  One common way to achieve this is to have a notion of
  \textbf{semantic value}, which differs from \textbf{syntactic
    value}. For example, the mathematical integers are part of the
  semantic values of Mini\hyp{}ML, but not of syntactical values.

\end{frame}

% ------------------------------------------------------------------------
%
\begin{frame}[fragile]{Semantic values}

  \begin{itemize}

    \item The \textbf{semantic values} are defined recursively by the
      following cases:\\
      \smallskip
      \begin{tabular}{@{\,}r@{\;\;}ll}
        $\bullet$
        & \emph{constant}
        & \unit, \textsf{0}, \textsf{1}, \textsf{2}, etc.\\
        $\bullet$
        & \emph{closure}
        & $\clos{x}{e}{\rho}$, with environment $\rho$.\\
        $\bullet$
        & \emph{$n$-tuple}
        & $v_1, \dots, v_n$, with $v_i$ values.
      \end{tabular}

    \item The closure environment $\rho$ is the one at the location
      where the function $\Xfun \; x \rightarrow e$ occurs, and the
      function is closed in that environment (this is called
      \textbf{static scoping}, or \textbf{lexical scoping}).

    \item The evaluation of an expression must take place in the
      context of an environment, where all free variables must be
      bound (in other words, the expression must be closed in that
      environment).

    \item For example, let $\rho$ be the environment before the
      sentence
\begin{alltt}
\small
\textbf{let} x = 1 \textbf{in} x
\end{alltt}
      Then the evaluation of the \emph{second} `x' must done in the
      environment \(\texttt{x} \mapsto 1 \oplus \rho\), where `x' is
      the \emph{first} one.

  \end{itemize}

\end{frame}

% ------------------------------------------------------------------------
%
\begin{frame}{Evaluation}

  Let us proceed to define a semantics for our mini-ML, by associating
  to each expression~\(e\) its value~\(v\), assuming the
  environment~\(\rho\):

  \bigskip

  \begin{tabular}{r@{\,\,}ll@{}}
    $\bullet$
  & $x$
  & Look up the \textbf{last} value bound to~\(x\) in~\(\rho\).\\
    $\bullet$
  & $\Xfun \; x \rightarrow e$
  & The value is the closure $\clos{x}{e}{\rho}$.\\
    $\bullet$
  & $e_1$ \texttt{+} $e_2$, etc.
  & Evaluate $e_1$ and $e_2$ in $\rho$, then add, etc.\\
    $\bullet$
  & \unit \ or \textsf{0} or \textsf{1} or \textsf{2}, etc.
  & The value is \unit or \textsf{0} or \textsf{1}, etc.\\
    $\bullet$
  & $\lpar{e}\rpar$
  & Evaluate $e$ into $v$ in $\rho$.\\
    $\bullet$
  & $\Xlet \; x = e_1 \Xin \; e_2$
  & Evaluate $e_1$ into $v_1$ in $\rho$;\\
  &
  & evaluate $e_2$ into $v$ in $x \mapsto v_1 \oplus \rho$.\\
    $\bullet$
  & $e_1 \; e_2$
  & Evaluate $e_1$ and $e_2$ into $v_1$ and $v_2$ in $\rho$\\
  &
  & (order left unspecified);\\
  &
  & $v_1$ must be the closure $\clos{x}{e}{\rho_1}$;\\
  &
  & evaluate $e$ in $x \mapsto v_2 \oplus \rho_1$:\\
  && the value is $v$.
\end{tabular}

\end{frame}

% ------------------------------------------------------------------------
%
\begin{frame}[fragile]{Example of evaluation}

  \begin{enumerate}

  \item The environment is initially empty. Let us evaluate
\begin{alltt}
\small
\textbf{let} x = 0 \textbf{in}
\textbf{let} id = \textbf{fun} x \(\rightarrow\) x \textbf{in}
\textbf{let} y = id x \textbf{in}
\textbf{let} x = (\textbf{fun} x \(\rightarrow\) \textbf{fun} y \(\rightarrow\) x + y) 1 2
\textbf{in} x + 1
\end{alltt}

  \item The evaluation of
\begin{alltt}
\small
\textbf{let} x = 0 \textbf{in} \ldots
\end{alltt}
    starts with the evaluation of~\texttt{\small 0}, which, obviously,
    yields~\(0\).

   \item We create the binding $\ident{x} \; \mapsto \; 0$, with which
     we extend~\(\rho_\varnothing\), that is, $\ident{x} \mapsto 0
     \oplus \rho_\varnothing$.

   \item We evaluate the elided subexpression within this new
     environment.

  \end{enumerate}

\end{frame}

% ------------------------------------------------------------------------
%
\begin{frame}{Example of evaluation}

  \begin{itemize}

    \item The evaluation of
\begin{alltt}
\small
\textbf{let} id = \textbf{fun} x \(\rightarrow\) x \textbf{in} \ldots
\end{alltt}
      is performed within the environment $\ident{x} \mapsto 0 \oplus
      \rho_\varnothing$.

    \item The value of~\texttt{id} is then the closure
      $\clos{\ident{x}}{\ident{x}}{\ident{x} \mapsto 0 \oplus \rho_\varnothing)}$.

    \item We extend the current environment and evaluate the elided
      subexpression within it.

  \item The evaluation of
\begin{alltt}
\small
\textbf{let} y = id x \textbf{in} \ldots
\end{alltt}
    is done within the environment $\rho = (\ident{id} \mapsto ((\Xfun
    \; \ident{x} \rightarrow \ident{x}), \ident{x} \mapsto 0 \oplus
    \rho_\varnothing)) \oplus (\ident{x} \mapsto 0) \oplus
    \rho_\varnothing$.

    \item We evaluate first (\texttt{\small id x}) in~\(\rho\).

    \item In order to do so, we evaluate~\texttt{\small id}
      and~\texttt{\small x}
      separately.

  \end{itemize}

\end{frame}

% ------------------------------------------------------------------------
%
\begin{frame}{Example of evaluation}

  \begin{itemize}

    \item These are both variables, thus we look them up in the
      environment to retrieve their last corresponding binding:
      \(\rho(\ident{id}) = \clos{\ident{x}}{\ident{x}}{\ident{x}
        \mapsto 0 \oplus \rho_\varnothing}\), and \(\rho(\ident{x}) =
      0\).

    \item We need to evaluate the body in the environment $\rho =
      (\ident{x} \mapsto 0) \oplus (\ident{x} \mapsto 0) \oplus
      \rho_\varnothing$, which yields~\(\rho(\ident{x}) = 0\), etc.

    \item We create the binding $\ident{y} \mapsto 0$, we add it to
      the current environment and we evaluate the elided subexpression
      with the new environment.

    \item Let us stop here and move to expressing the above semantics
      in a formal manner.

  \end{itemize}

\end{frame}

% ------------------------------------------------------------------------
%
\begin{frame}{Formal semantics}

  If we note \(\evalf{e}{\rho}\) the value obtained by evaluating the
  expression~\(e\) in the environment~\(\rho\), then the evaluation of
  mini-ML is:
  \begin{align*}
    \evalf{\overline{n}}{\rho} &= n, \;\, \text{with
      $\overline{n}$ an \OCaml integer and $n \in \mathbb{N}$}\\
    \evalf{e_1 \; \texttt{+} \; e_2}{\rho} &=
    \evalf{e_1}{\rho} + \evalf{e_2}{\rho} \in \mathbb{N}\\
    \evalf{\lpar{e}\rpar}{\rho} &= \evalf{e}{\rho}\\
    \evalf{x}{\rho} &= \rho(x)\\
    \evalf{\Xfun \; x \rightarrow e}{\rho} &= \clos{x}{e}{\rho}\\
    \evalf{e_1 \; e_2}{\rho} &= \evalf{e_3}{(x \mapsto
      \evalf{e_2}{\rho} \oplus \rho_3)},\\
    &\phantom{=\llbracket} \textnormal{where} \ \evalf{e_1}{\rho} =
    \clos{x}{e_3}{\rho_3}\\
    \evalf{\Xlet \; x = e_1 \; \Xin \; e_2}{\rho}
    &= \evalf{(\Xfun \; x \rightarrow e_2) \; e_1}{\rho}.
  \end{align*}

  Note that we write ``$x,\, y$'' instead of ``$(x,\, y)$'' for
  brevity. Evaluation means applying the equations from left to right.

\end{frame}

% ------------------------------------------------------------------------
%
\begin{frame}[fragile]{Partial applications}

  \begin{itemize}

    \item We can easily work out that
      \begin{equation*}
        \evalf{\Xlet \; x = e_1 \; \Xin \; e_2}{\rho}
        = \evalf{e_2}{(x \mapsto \evalf{e_1}{\rho} \oplus \rho)}.
      \end{equation*}

    \item Since closures are values, they can be the value of a
      function call:
\begin{alltt}
\small
\textbf{let} add = \textbf{fun} x \(\rightarrow\) \textbf{fun} y \(\rightarrow\) x + y \textbf{in}
\textbf{let} incr = add 1
\textbf{in} incr 5
\end{alltt}

    \item The call (\texttt{\small add 1}) is a \textbf{partial
      application}, as opposed to a \textbf{complete application} like
      (\texttt{\small add 1 5}). A partial application, by definition,
      always evaluates in a closure.

    \item Operations are functions denoted by a symbol used in infix
      position (between the operands). As such, they can be partially
      evaluated:
\begin{alltt}
\small
\textbf{let} incr = (+) 1 \textbf{in} incr 5
\end{alltt}
     \smallskip
     \noindent where the parentheses in~\texttt{(+)} mean that the
     operator has to be considered in \textbf{prefix} position.

  \end{itemize}

\end{frame}

% ------------------------------------------------------------------------
%
\begin{frame}[fragile]{Auto-application and termination}

  \begin{itemize}

    \item Our mini-ML is \textbf{Turing\hyp{}complete}, just as \OCaml
      is.

    \item For instance, we can write the following
      non\hyp{}terminating program:
      \smallskip
\begin{alltt}
\small
\textbf{let} omega = \textbf{fun} f \(\rightarrow\) f f \textbf{in} omega omega
\end{alltt}

    \item The function \texttt{\small omega} is an
      \textbf{auto\hyp{}applicative function}.

    \item It enables the following non\hyp{}terminating evaluation:
      \begin{gather*}
        \evalf{\texttt{\textbf{let} omega = \textbf{fun} f \(\rightarrow\) f f \textbf{in} omega omega}}{\rho}\\
\begin{align*}
&= \evalf{\texttt{\small omega omega}}{(\texttt{\small omega} \mapsto
    \evalf{\texttt{\small \textbf{fun} f \(\rightarrow\) f f}}{\rho} \oplus
    \rho)}\\
&= \evalf{\texttt{\small f f}}{(\texttt{\small f} \mapsto \clos{f}{f \; f}{\rho})
    \oplus \rho)}\\ &= \text{\emph{idem.}}
\end{align*}
\end{gather*}

  \end{itemize}

\end{frame}

% ------------------------------------------------------------------------
%
\begin{frame}[fragile]{Fixed points}

  \begin{itemize}

    \item To demonstrate the expressive power of our mini-ML, let us
      see how we can simulate recursive functions by means of the
      $\Omega$ function:
\begin{alltt}
\small
\textbf{let} \(\Omega\) = \textbf{fun} f \(\rightarrow\) f f
\end{alltt}

    \item Let us define a function \texttt{\small fix}, called the
      \textbf{Y combinator} in \(\lambda\)-calculus:
      \smallskip
\begin{alltt}
\small
\textbf{let} fix = \textbf{fun} g \(\rightarrow\) \(\Omega\) (\textbf{fun} h \(\rightarrow\) \textbf{fun} x \(\rightarrow\) g (h h) x)
\end{alltt}

      \item For any (terminating) function~\(f\) and
        any function~\(x\),
        \begin{equation*}
          \evalf{f \; \lpar\ident{fix} \; f\rpar{} \; x}{\rho} =
          \evalf{\ident{fix} \; f \; x}{\rho}.
        \end{equation*}

      \item In other words, we have
        \begin{equation*}
          f \; (\ident{fix} \; f) \; x \equiv \ident{fix} \; f \; x.
        \end{equation*}

      \item The fixed point~\(x\) of a (continuous) function~\(g\)
        satisfies \(g(x) = x\).

      \item Therefore, the fixed point of~\ident{f} (if it exists) is
        the value of \lpar\ident{fix} \ident{f}\rpar{}.

  \end{itemize}

\end{frame}

% ------------------------------------------------------------------------
%
\begin{frame}{Fixed points}

   Just as with \Xlet \Xin, we can simplify the semantics of
   \texttt{\small fix}:
   \begin{equation*}
     \evalf{\ident{fix} \; e}{\rho} = \evalf{e_1}{(f \mapsto
       \evalf{\ident{fix} \, (\Xfun \, f \rightarrow e_1)}{\rho}
       \oplus \rho_1)},
   \end{equation*}
   where $\evalf{e}{\rho} = \clos{f}{e_1}{\rho_1}$.

   \bigskip

   There is a choice:
   \begin{itemize}

     \item either we define \texttt{\small fix} in the functional
       language (as a built\hyp{}in constant), or

     \item at the meta\hyp{}language (the mathematical logic we are
       using to define the semantics).

   \end{itemize}

\end{frame}

% ------------------------------------------------------------------------
%
\begin{frame}[fragile]{Fixed points}

  Let us consider the following definitions:
\begin{alltt}
\small
\textbf{let} pre\_fact =
  \textbf{fun} f \(\rightarrow\) \textbf{fun} n \(\rightarrow\) \textbf{if} n = 1 \textbf{then} 1 \textbf{else} n * f(n-1)
\textbf{let} fact = fix pre\_fact
\end{alltt}

  We see that \texttt{\small fact} is the fixed point of
  \texttt{\small pre\_fact}:
  \begin{align*}
    \evalf{\texttt{\small fact}}{\rho} &= \evalf{\texttt{\small fix
        pre\_fact}}{\rho}\\
    &= \evalf{\texttt{\small pre\_fact (fix pre\_fact)}}{\rho}\\
    &= \evalf{\texttt{\small pre\_fact fact}}{\rho}\\
    &= \evalf{\texttt{\small \textbf{fun} n \(\rightarrow\) \textbf{if} n
        = 1 \textbf{then} 1 \textbf{else} n * fact(n-1)}}{\rho}
  \end{align*}

  It is as if the definition of `\texttt{\small fact}' were
  recursive... but it is not.

\end{frame}

% ------------------------------------------------------------------------
%
\begin{frame}[fragile]{Fixed points in C}

\begin{alltt}
\small
\textbf{#include}<stdio.h>
\textbf{#include}<stdlib.h>

\textbf{typedef int} (*fp)();

\textbf{int} fact(fp f, \textbf{int} n) \{
  \textbf{return} n? n * ((\textbf{int} (*)(fp,\textbf{int}))f)(f,n-1) : 1; \}

\textbf{int} read(\textbf{int} dec, \textbf{char} arg[]) \{
  \textbf{return} ('0' <= *arg && *arg <= '9')?
         read(10*dec + *arg - '0', arg+1) : dec; \}

\textbf{int main}(\textbf{int} argc, \textbf{char}** argv) \{
  \textbf{if} (argc == 2)
     printf("%u\textbackslash{n}", fact(&fact, read(0, argv[1])));
  \textbf{else} printf("Only one integer allowed.\textbackslash{n}");
  \textbf{return} 0; \}
\end{alltt}
\end{frame}

%% % ------------------------------------------------------------------------
%% %
%% \begin{frame}{Local recursive definitions}

%%   \begin{itemize}

%%     \item The function `\texttt{\small fact}' coincides pointwise with
%%       the factorial function.

%%     \item Let us extend mini-ML with a native recursive binder, using
%%       the `\texttt{\small fix}' function:
%%       \begin{equation*}
%%         \evalf{\Xlet \; \Xrec \; f = e_1 \; \Xin \; e_2}{\rho} =
%%         \evalf{\Xlet \; f = \textsf{fix} \; (\Xfun \; f \rightarrow
%%           e_1) \; \Xin \; e_2}{\rho}.
%%       \end{equation*}

%%   \end{itemize}

%% \end{frame}

%% % ------------------------------------------------------------------------
%% %
%% \begin{frame}{Extensions}

%%   \begin{itemize}

%%     \item Let us extend our mini-ML with the following
%%       expressions.\\
%%       \smallskip
%%       \begin{tabular}{r@{\,\,}ll@{}}
%%         $\bullet$
%%         & \emph{Boolean constants}
%%         & \Xtrue{} or \Xfalse\\
%%         $\bullet$
%%         & \emph{Boolean operators}
%%         & \texttt{(\&\&)} or \texttt{(||)} or \textsf{\textbf{not}}\\
%%         $\bullet$
%%         & \emph{$n$-tuple}
%%         & $e_1, \ldots, e_n$\\
%%         $\bullet$
%%         & \emph{conditional}
%%         & \Xif{} $e_0$ \Xthen{} $e_1$ \Xelse{} $e_2$\\
%%         $\bullet$
%%         & \emph{local recursive binding}
%%         & \Xlet{} \Xrec{} $f \;\; = \;\; e_1$ \Xin{} $e_2$
%%       \end{tabular}

%%       \item Note that parentheses are recommended around tuples, but
%%         are not mandatory.

%%       \item Let us distinguish the variables occurring after \Xlet{}
%%         and \Xfun{}, because they are \textbf{irrefutable patterns},
%%         noted $\ipat{p}$:\\
%%         \smallskip
%%         \begin{tabular}{r@{\,\,}ll@{}}
%%           $\bullet$
%%           & \emph{function}
%%           & $\Xfun \;\; \ipat{p} \rightarrow e$\\
%%           $\bullet$
%%           & \emph{local definition}
%%           & \(\Xlet \;\; \textrm{[}\Xrec\textrm{]} \;\; \ipat{p} \;\;
%%           = \;\; e_1 \;\; \Xin \;\; e_2\)
%%         \end{tabular}

%%   \end{itemize}

%% \end{frame}

%% % ------------------------------------------------------------------------
%% %
%% \begin{frame}{Irrefutable patterns}

%%   \begin{itemize}

%%     \item An irrefutable pattern $\ipat{p}$ is defined recursively by
%%       the following cases:\\
%%       \smallskip
%%       \begin{tabular}{r@{\,\,}ll@{}}
%%         $\bullet$
%%         & \emph{variable}
%%         & $f, g, h, x, y, z$, etc. \\
%%         $\bullet$
%%         & \emph{unit}
%%         & \unit\\
%%         $\bullet$
%%         & \emph{$n$-tuple}
%%         & $\ipat{p}_1, \ldots, \ipat{p}_n$\\
%%         $\bullet$
%%         & \emph{parenthesis}
%%         & $\lpar\ipat{p}\rpar$\\
%%         $\bullet$
%%         & \emph{underscore}
%%         & {\Large \_}
%%       \end{tabular}

%%       \item Note that, from the syntactic point of view, irrefutable
%%         patterns are special cases of expressions, except the
%%         underscore, which is a special case avoiding the creation of a
%%         binding.

%%       \item The sentence~\(e\) is equivalent to \(\Xlet \;\; \_ \;\; =
%%         \;\; e\).

%%   \end{itemize}

%% \end{frame}


\end{document}
